% !TEX root = ../thesis-main.tex

\paperchapter[Barcode matchings and $p$-norms]{GRRS2026}
\newrefsection

% Authors and affiliations
\setcounter{footnotevalue}{\value{footnote}}
\setcounter{footnote}{0}
\renewcommand*{\thefootnote}{\roman{footnote}}
{\sffamily
  Andrea Guidolin\footnote{University of Southampton},
  Arthur Reidenbach\footnote{ENS Paris--Saclay},
  Isaac Ren\footnote{KTH Royal Institute of Technology},
  Martina Scolamiero\addtocounter{footnote}{-1}\footnotemark
}
\setcounter{footnote}{\value{footnotevalue}}
\renewcommand*{\thefootnote}{\arabic{footnote}}
\vspace{1ex}

\begin{paperabstract}
  We construct barcode matchings induced by morphisms of pointwise finite-dimensional persistence modules. To do this, we use an algebraic decomposition of the space of morphisms between persistence modules.
  Applying our induced matchings to arbitrary morphisms by the epi-mono factorization strategy of Bauer and Lesnick (2015), we recover the partial additive matching of Gonzalez-Diaz, Sorriano-Trigueros, and Torras-Casas (2025).
  We also study the interaction between our matching and a notion of cost for morphisms based on $p$-norms of persistence modules. We illustrate some consequences of our results on this cost in the context of algebraic stability of persistence barcodes, including a simple proof of the algebraic Wasserstein isometry of persistent homology.
\end{paperabstract}

\section{Introduction}

One of the central results in topological data analysis (TDA) is the barcode decomposition theorem, which states that, under some mild assumptions, every one-dimensional persistence module is isomorphic to a direct sum of basic interval modules, parameterized by intervals of the indexing poset. We can then represent the algebraic data of a persistence module by the combinatorial multiset of intervals. In this paper, we consider \textbf{pointwise finite-dimensional} (\textbf{p.f.d.}) \textbf{persistence modules} over a fixed field $K$, which are defined as the functors from the poset $(\R, \mathord{\leq})$ to the category $\Vect_K$ of finite dimensional vector spaces. In particular, p.f.d.\ modules admit a barcode decomposition theorem \cite{CrawleyBoevey2015,BotnanCrawleyBoevey2020}.

In this context, it is natural to ask whether morphisms of persistence modules also admit a combinatorial representation. The natural candidate for such representatives are (partial) \textbf{matchings}, or partial bijections, between barcodes. Barcode matchings are in fact quite central to the stability theory of persistent homology, as they form the basis of the definitions of the bottleneck and Wasserstein distances \cite[Section VIII.2]{EdelsbrunnerHarer2010}. Barcode matchings induced by morphisms of persistence modules were studied previously by Bauer and Lesnick \cite{BauerLesnick2015,BauerLesnick2020}, who recover various stability results for p.f.d.\ modules, by Skraba and Turner \cite{SkrabaTurner2023}, who use induced matchings to show an isometry result between algebraic and combinatorial Wasserstein distances, and by Gonzalez-Diaz, Sorriano-Trigueros, and Torras-Casas \cite{gonzalez2021partial,GST2025}, who construct an additive barcode matching for persistence modules with finite barcodes.

Going back to the question of combinatorial representations of morphisms, Jacquard, Nanda, and Tillman provide a decomposition theorem for morphisms between persistence modules with finite barcodes that do not contain strictly nested bars \cite[Theorem 5.3]{JNT2023}. More precisely, given a morphism of persistence modules $f \colon \bigoplus_{i \in I} K(a_i, b_i) \to \bigoplus_{j \in J} K(c_j, d_j)$, where $K(a, b)$ denote the interval module with support $[a, b)$, if there is no pair of indices $i, i'$ in $I$ such that $[a_i, b_i)$ is a strict subset of $[a_{i'}, b_{i'})$, and no similar pair $j, j'$ in $J$, then there exists a partial bijection $\mu \subseteq I \times J$ (i.e., each element of $I$ and $J$ appears at most once in $\mu$) such that
\begin{equation}
\label{C!E:intro-btb}
  f \cong \bigoplus_{(i, j) \in \mu} f_{i, j} \oplus \bigoplus_{i \in I \setminus \pi_I(\mu)} g_i \oplus \bigoplus_{j \in J \setminus \pi_J(\mu)} h_j , 
\end{equation}
where each $f_{ij}$ is a nonzero morphism $K(a_i, b_i) \to K(c_j, d_j)$, each $g_i$ denotes the zero morphism $K(a_i, b_i) \to 0$ and each $h_j$ denotes the zero morphism $0 \to K(c_j, d_j)$. In \cite{AGRS2025}, the authors (some of whom are authors of the current paper) call a morphism \textbf{bar-to-bar} if it admits a decomposition as in \eqref{C!E:intro-btb}. Such bar-to-bar morphisms can be seen as an algebraic version of barcode matchings; indeed, $\mu$ is the corresponding matching. In particular, the authors provide an algorithm for inducing bar-to-bar morphisms from mono- and epimorphisms.

\subsection{Our contributions}

In this paper, we continue the analysis of \cite{AGRS2025} by studying the association of matchings to morphisms. Following the strategy of \cite{BauerLesnick2015}, we induce a barcode matching for an arbitrary morphism $f \colon X \to Y$ by first finding an epi-mono factorization
\[
  f \colon X \xrightarrow{e} \im(f) \xrightarrow{m} Y
\]
where $\im(f)$ denotes the image of $f$, $e \colon X \to \im(f)$ an epimorphism, and $m \colon \im(f) \to Y$ a monomorphism. We then compute the matchings of $e$ and $m$ associated to their induced bar-to-bar morphisms, and then compose them together to obtain a matching associated with $f$. Note that this depends \textit{a priori} on the choice of epi-mono factorization; in practice, this is not the case for the induced matching of \cite{BauerLesnick2015} but is the case for our matchings (\cref{C!ex:counterex2}).

The key insight of this paper is that we can induce bar-to-bar morphisms by projecting morphisms down to a subspace of morphisms that only map nontrivially between bars with the same end-point: this is the \textbf{same-endpoint operator}, denoted by $\Delta$ (\cref{C!D:Delta}). Given a morphism $f$, $\Delta f$ is a direct sum of morphisms between persistence modules with no nested bars. We prove, in the p.f.d.\ setting, that these morphisms are bar-to-bar, which implies that $\Delta f$ is bar-to-bar (\cref{C!T:Delta-is-btb}). Notably, this simplifies the construction of induced bar-to-bar morphisms from \cite{AGRS2025}. As described above, this naturally induces a barcode matching $\mu_\Delta(f)$ (\cref{C!D:Delta-matching}). In particular, this matching is additive: given morphisms $f$ and $g$, $\mu_\Delta(f \oplus g) = \mu_\Delta(f) \oplus \mu_\Delta(g)$.

Our proof techniques rely on matrix representations of morphisms. To this end, we introduce the notion of labeled matrices and their multiplication (\cref{C!D:labeled-matrix}). In particular, to show that $\Delta f$ is bar-to-bar, we reduce a matrix representation of $\Delta f$ to a matrix with at most one nonzero coefficient in each row and column (\cref{C!A:inf-red}). The induced matching is then read from the pivots of this reduced matrix representation. This matrix reduction is similar to \cite[Theorem 5.3]{JNT2023}, the decomposition theorem for morphisms of persistence modules without nested bars.

\subsection{Organization of the paper}

In \cref{C!S:preliminaries}, we define the notions used in this paper, including p.f.d.\ persistence modules, labeled matrices, barcode matchings, and dualization results.

In \cref{C!S:btb}, we introduce the notion of bar-to-bar morphisms and describe how to induce barcode matchings from an arbitrary morphisms using the same-end operator and epi-mono factorizations. In particular, we recover the $\cc{P}$-matching of \cite{GST2025} in \cref{C!SS:P-matching}. Along the way, we obtain a simple algorithm for computing the barcode of the image of a morphism of persistence modules (\cref{C!P:image-barcode}).

We apply bar-to-bar morphisms to (algebraic) Wasserstein distances in \cref{C!S:Wasserstein}. In \cref{C!SS:coker}, we provide a simpler proof of \cite[Theorem 3.14]{AGRS2025} (while also generalizing the result to p.f.d.\ modules), which states that, for a given monomorphism $f \colon X \hookrightarrow Y$, there exists a bar-to-bar monomorphism $f_b \colon X \hookrightarrow Y$ such that $\| \coker f_b \|_p \leq \| \coker f \|_p$, where $\| - \|_p$ denotes the $p$-norm of p.f.d.\ persistence modules (\cref{C!thm:cost-coker-btb}). As part of this new proof, we obtain an expression of the $p$-norm of a persistence module as an integral of the rank function (\cref{C!P:p-norm-by-integral}). In \cref{C!SS:isometry}, we use the $p$-norm inequality result, and the dual result for epimorphisms, to prove the isometry between the algebraically defined and combinatorially defined Wasserstein distances (\cref{C!T:isometry}).

\section{Preliminaries}
\label{C!S:preliminaries}

\subsection{Persistence modules} 

We fix a field $K$ and denote by $\Vect_\mathit{K}$ the category of finite dimensional vector spaces over $K$. Viewing $(\R, \mathord{\leq})$ as a poset category, a \index{persistence module!pointwise finite-dimensional (p.f.d.)}\textbf{pointwise finite-dimensional persistence module over $K$} is a functor $X \colon \R \to \Vect_\mathit{K}$, consisting therefore of a collection of finite dimensional vector spaces $\{X_t\}_{t \in \mathbb{R}}$ and linear maps between them $X_{s \leq t}$ for $s \leq t$,  we call transition functions. We denote the category of pointwise finite-dimensional persistence modules and natural transformations between them by \index{.b pfd@$\Pfd$}$\Pfd$. In the following, pointwise finite-dimensional persistence modules are referred to as persistence modules for brevity. The category $\Pfd$ has biproducts (i.e.\ direct sums), which are computed pointwise\footnote{It is a general fact that limits and colimits in functor categories are computed pointwise; $\Pfd$ is a functor category and biproducts are both limits and colimits.}.

In this setting, interval modules are not parameterized by the real line, but rather by an extension of the real line:

\begin{definition}[Decorated real line and intervals]
\label{C!def:decoratedE}
  This definition is adapted from \cite[\textsection 2.3]{BauerLesnick2015}.
  We consider the \textbf{decorated real line}
  \[
    \E \coloneq \R \times \{-, +\} \cup \{(-\infty, +), (+\infty, -)\} \subseteq [-\infty, +\infty] \times \{-, +\}
  \]
  equipped with the following lexicographic order: for $(s, \sigma)$ and $(t, \tau)$ in $\E$, we have $(s, \sigma) < (t, \tau)$ if $s < t$ in $[-\infty, +\infty]$ or if $s = t$, $\sigma = -$, and $\tau = +$. We denote $(t, -)$ and $(t, +)$ in $\E$ by $t^-$ and $t^+$, respectively, except when $t = \pm \infty$, where we omit the sign.

  Ordered pairs of elements of $\E$ are in one-to-one correspondence with all intervals on the real line: for $s^\sigma < t^\tau$ in $\E$, we define the interval
  \[
    \langle s^\sigma, t^\tau \rangle \coloneq \{u \in \R \mid s < u < t \text{ or } s^\sigma = u^- \text{ or } t^\tau = u^+\}.
    \qedhere
  \]
\end{definition}

In particular, for $s < t$ in $\R$, we have the following equality of intervals:
\begin{align*}
  \langle s^-, t^- \rangle & = [s, t), &
  \langle s^-, t^+ \rangle & = [s, t], \\
  \langle s^+, t^- \rangle & = (s, t), &
  \langle s^+, t^+ \rangle & = (s, t].
\end{align*}

\begin{definition}[Overlapping intervals]
  Let $a, b, c, d \in \E$. Following \cite{BauerLesnick2020}, we say that \index{overlap above/below}\textbf{the interval $\langle a, b \rangle$ overlaps the interval $\langle c, d \rangle$ above} (and antisymmetrically, $\langle c, d \rangle$ \textbf{overlaps $\langle a, b \rangle$ below}) if $c \leq a < d \leq b$.
\end{definition}

\begin{definition}[Interval modules]
  Let $a < b$ be elements of $\E$.
  A \textbf{bar} (or \textbf{interval module}) with \index{start-point}\textbf{start-point} $a$ and \index{end-point}\textbf{end-point} $b$, denoted as $K(a, b)$, is the persistence module defined as follows: for any $u$ in $\R$,
  \begin{equation*}
    K(a, b)_u \coloneq
    \begin{cases}
      K & \text{if } u \in \langle a, b \rangle, \\
      0 & \text{otherwise},
    \end{cases}
  \end{equation*}
  and for any $u \leq v$ in $\R$,
  \begin{equation*}
    K(a, b)_{u \leq v} \coloneq
    \begin{cases}
      \id_K & \text{if } u, v \in \langle a, b \rangle, \\
      0 &\text{otherwise}.
    \end{cases}
  \end{equation*}
  In particular, the \textbf{support} of $K(a, b)$, i.e.\ the real values where the functor is nonzero, is the interval $\langle a, b \rangle$.
\end{definition}

\begin{remark}
\label{C!R:overlap}
  We observe that there exists a nonzero morphism $K(a, b) \to K(c, d)$ if, and only if, $\langle a, b \rangle$ overlaps $\langle c, d \rangle$ above.
\end{remark}

In order to better manage these interval modules, we will be considering the following subcategory of $\Pfd$. Recall that a \index{multiset}\textbf{multiset} is a set $S$ equipped with a \index{multiplicity (for multisets)}\textbf{multiplicity} $m_S \colon S \to \N$, and that when there is no ambiguity, $S$ and $m_S$ may be used interchangeably.

\begin{definition}[Subcategory of bars]
  We denote by \index{.b bar@$\Bar$}$\Bar$ the full subcategory of $\Pfd$ whose objects are the (potentially infinite) direct sums of bars $\bigoplus_{i \in I} K(a_i, b_i)$. We call this the \textbf{category of bars}.

  Given an object $X = \bigoplus_{i \in I} K(a_i, b_i)$ of $\Bar$, we define its \index{barcode}\textbf{barcode} as a representation of the multiset $\cc{B}(X) \coloneq \{\langle a_i, b_i \rangle\}_{i \in I}$ of intervals appearing in $X$ (recall the definition of multiset representations from \cite{BauerLesnick2020}).
\end{definition}

\begin{definition}
\label{C!D:barcode-decomp}
  Let $X$ be a persistence module in $\Pfd$. A \index{barcode decomposition}\textbf{barcode decomposition} of $X$ is the given of an object $\bigoplus_i K(a_i, b_i)$ in $\Bar$ and an isomorphism $\varphi_X \colon X \cong \bigoplus_i K(a_i, b_i)$.
\end{definition}

\begin{remark}
  The decomposition theorem \cite{CrawleyBoevey2015,BotnanCrawleyBoevey2020} states that every persistence module  $X$ in $\Pfd$ admits a barcode decomposition  $\varphi_X \colon X \cong \bigoplus_i K(a_i, b_i)$. While the isomorphism $\varphi_X$ is not unique, the direct sum of bars $\bigoplus_i K(a_i, b_i)$ is, up to reordering of the bars. In particular, we can speak of the \textbf{barcode} of an object in $\Pfd$ as well.  
\end{remark}

\begin{remark}
\label{C!R:ephemeral}
  Every persistence module in $\Pfd$ has at most countably many bars $K(s, t)$ with $s < t \in E$ in its barcode decomposition. In other words, up to an \index{ephemeral module}\textbf{ephemeral} summand (see \cite{CdSGO2016}), which is a direct sum of bars of the form $K(s^-, s^+)$ supported on $\{s\}$, every persistence module $X$ in $\Pfd$ has at most countably many bars. 
  Moreover, for each $s \in \R$, $X$ has a finite number of bars $K(s^-, s^+)$, since the dimension of $X_s$ is finite. As a consequence, for all $b \in \E$, $X$ only has countably many bars with start-point $b$, and likewise countably many with end-point $b$.
\end{remark}


\begin{proposition}
  The categories $\Bar$ and $\Pfd$ are equivalent, given the axiom of choice.
\end{proposition}

\begin{proof}
  We observe that the inclusion functor $\Bar \hookrightarrow \Pfd$ is fully faithful as the inclusion of a full subcategory and is essentially surjective by the decomposition theorem. Using the axiom of choice, we conclude that the inclusion functor is an equivalence of categories: see \cite[\href{C!https://stacks.math.columbia.edu/tag/02C3}{Lemma 02C3}]{stacks-project}.
\end{proof}

Recall from \cite{AGRS2025} that a persistence module $X$ in Pfd is \textbf{tame} if there exist real numbers $0=t_0<t_1<\cdots <t_k$ such that the transition function $X_{s\le t}$ is a non-isomorphism only if $s<t_i \le t$ for some $i\in \{ 1,\ldots ,k\}$.  
All the bars in the barcode decompositon of a tame persistence module are of the form $\langle s^-, t^- \rangle = [s, t)$.

As a notion of \emph{size} for persistence modules, we use the $p$-norm first introduced in \cite{p-norm_Edesbrunner} and \cite{SkrabaTurner2023}, also studied in \cite{AGRS2025} for tame persistence modules.  

\begin{definition}
\label{C!def:p-norm}
  For $p\in [1,\infty]$, the \index{.a p-norm@$p$-norm}\textbf{$p$-norm} of a persistence module with barcode decomposition $X\cong \bigoplus_{i} K(a_i, b_i)$ is:
  \begin{equation*}
    \left\| X\right\|_{p} =  \begin{cases}
      \left( \sum_{i} (b_i-a_i)^{p} \right)^{\frac{1}{p}} &\text{for $p\in [1,\infty)$},\\
      \sup_i \{ (b_i-a_i) \} &\text{for $p=\infty$} .
    \end{cases}
    \qedhere
  \end{equation*}
\end{definition}

By Remark \ref{C!R:ephemeral}, the $p$-norm of a persistence module in $\Pfd$ involves an infinite sum with at most countably many positive summands. Following \cite{SkrabaTurner2023} (Definition $2.11$) we say that a persistence module $X$ has \emph{bounded p-norm} if $\left\| X\right\|_{p}^p<\infty$.

\subsection{Morphisms and matrices}
\label{C!sec:mor-mtrx}

One advantage of $\Bar$ over $\Pfd$ is that we can conveniently, and not just up to isomorphism, represent morphisms as (possibly infinite) matrices. Let $X = \bigoplus_{i \in I} K(a_i, b_i)$ and $Y = \bigoplus_{j \in J} K(c_j, d_j)$ be objects of $\Bar$. Note that $\hom(X, Y)$ is a vector space, possibly infinite-dimensional. Moreover, given the canonical isomorphisms
\begin{equation}
\label{C!E:iso1}
  \hom (\bigoplus_{i \in I} K(a_i, b_i), \bigoplus_{j \in J} K(c_j, d_j)) \cong \prod_{i \in I} \prod_{j \in J}  \hom (K(a_i, b_i), K(c_j, d_j))
\end{equation}
and
\begin{equation}
\label{C!E:iso2}
  \hom ( K(a_i, b_i),  K(c_j, d_j)) \cong 
  \begin{cases}
    K & \text{if } \langle a_i, b_i \rangle \text{ overlaps } \langle c_j, d_j \rangle \text{ above}, \\
    0 & \text{otherwise},
  \end{cases}
\end{equation}
we can then represent each morphism $f \in \hom(X, Y)$ with a matrix $F \in K^{I \times J}$, where $F_{ij} = 0$ if $\langle a_i, b_i \rangle$ does not overlap $\langle c_j, d_j \rangle$ above.

\begin{definition}[Column-finite matrices]
\label{C!D:column-finite-matrices}
  A matrix $M \in K^{I \times J}$ with rows indexed by $I$ and columns indexed by $J$ is \textbf{column-finite} if each column only has a finite number of nonzero coefficients.
\end{definition}

\begin{definition}[Labeled matrices]
\label{C!D:labeled-matrix}
  A \index{labeled matrix}\textbf{labeled matrix} is a (potentially infinite) matrix $A \in K^{I \times J}$ where $I$ and $J$ are the totally ordered indexing sets for rows and columns, respectively, and where each row and column is labeled by an interval of $\R$ such that, for all $(i, j) \in I \times J$, writing $\langle a_i, b_i \rangle$ for the row label and $\langle c_j, d_j \rangle$ for the column label, $A_{ij}$ is nonzero only if $\langle c_j, d_j \rangle$ overlaps $\langle a_i, b_i \rangle$ above.

  Let $A \in K^{I \times J}$ and $B \in K^{J \times L}$ be labeled, column-finite matrices where the co\-lumns of $A$ and the rows of $B$ have the same labels. The \index{overlap multiplication}\textbf{overlap multiplication} of $A$ and $B$, written as the usual product $AB$, is defined by, for all $(i, \ell) \in I \times L$,
  \begin{equation}
  \label{C!E:labeled-matrix}
    (AB)_{i \ell} = \sum_{j \in J} A_{ij} B_{j\ell} \bb{I}\{a_i \leq e_\ell < b_i \leq f_\ell\},
  \end{equation}
  where the label of row $i$ in $A$ is $\langle a_i, b_i \rangle$ and the label of column $\ell$ in $B$ is $\langle e_\ell, f_\ell \rangle$, and $\bb{I}\{a_i \leq e_\ell < b_i \leq f_\ell\} = 1$ if the inequalities hold and $0$ otherwise.
\end{definition}

\begin{remark}
  \begin{enumerate}
    \item We check that overlap multiplication is associative, with the usual identity matrix as the unit.
    \item The reason we assume that $F$ and $G$ are column-finite is to ensure that the sums in \eqref{C!E:labeled-matrix} are well-defined. There are other settings where the sums are well-defined, e.g.\ for row-finite matrices, but we will only need the result for column-finite matrices.
    \item When every interval indexed by $L$ overlaps every interval indexed by $I$, overlap multiplication corresponds exactly to standard matrix multiplication. For example, this is the case when every interval has the same end-point. \qedhere
  \end{enumerate}
\end{remark}

\begin{definition}[Matrix representations of bar morphisms]
\label{C!D:matrix}
  Let $f \colon \bigoplus_{j \in J} K(a_j, b_j) \to \bigoplus_{i \in I} K(c_i, d_i)$ be a morphism in $\Bar$ (note the swapped indices). An \index{matrix representation!ordered}\textbf{ordered matrix representation of $f$} is a labeled matrix $(F_{ij})_{i \in I, j \in J}$ where
  \begin{enumerate}
    \item for all $i \in I$ and $j \in J$, the entry $F_{ij}$ is the projection of the morphism $f$ to $\hom(K(a_j, b_j), K(c_i, d_i))$ via the canonical isomorphisms \eqref{C!E:iso1} and \eqref{C!E:iso2},
    \item $I$ and $J$ are ordered in a way compatible with the product order on the pairs $(a_j, b_j)$ and $(c_i, d_i)$, and
    \item row $i$ is labeled by $\langle c_i, d_i \rangle$ and column $j$ is labeled by $\langle a_j, b_j \rangle$. \qedhere
  \end{enumerate}
\end{definition}

\begin{example}
  For example, the only ordered matrix representation of the morphism represented in Figure \ref{C!fig:morphism_in_Bar} is:
  \[
    F = \begin{blockarray}{c cc}
      & \langle a_1,b_1 \rangle & \langle a_2,b_2 \rangle \\
      \begin{block}{c [cc]}
        \langle c_1,d_1 \rangle & 1 & 1 \\
        \langle c_2,d_2 \rangle & 1 & 0 \\
      \end{block}
    \end{blockarray},
    \qedhere
  \]
\end{example}

\begin{figure}[tbp]
  \centering
    \begin{tikzpicture}[xscale=1.5, yscale=1, thick]
      \draw[-] (1.8,2.5) node[left] {$a_1$} -- (2.8,2.5) node[right] {$b_1$};
      \draw[-] (1.4,1.8) node[left] {$a_2$} -- (3.2,1.8) node[right] {$b_2$};

      \draw[-] (0.7,0.5) node[left] {$c_1$} -- (2.2,0.5) node[right] {$d_1$};
      \draw[-] (0.8,-0.2) node[left] {$c_2$} -- (2.8,-0.2) node[right] {$d_2$};

      \draw[->, thin] (2,2.4) -- (2,0.6);
      \draw[->, thin] (1.5,1.7) -- (1.5,0.6);
      \draw[->, thin] (2.6,2.4) -- (2.6,-0.1);

      \node (X) at (4.5, 2.3) {$X$};
      \node (Y) at (4.5, 0) {$Y$};
      \draw[->] (X) -- node[right] {$f$} (Y);
    \end{tikzpicture}
  \caption{Representation of a morphism $f \colon K(a_1,b_1) \oplus K(a_2,b_2) \rightarrow K(c_1,d_1) \oplus K(c_2,d_2)$ in $\Bar$. An arrow between two bars $K(a_i,b_i)$ and $K(c_j,d_j)$ corresponds to the projection of $f$ to $\hom(K(a_i, b_i), K(c_j, d_j))$; the morphism is represented by its nonzero projections.}
  \label{C!fig:morphism_in_Bar}
\end{figure}

Under certain conditions, the composition of morphisms can also be represented by matrices.

\begin{proposition}
\label{C!P:matrix-multiplication}
  Let $f \colon X \to Y$ and $g \colon Y \to Z$ be morphisms in $\Bar$, and $F$ and $G$ matrix representations of theirs. Suppose that $F$ and $G$ are column-finite. Then the overlap multiplication $GF$ (\cref{C!D:labeled-matrix}) is a matrix representation of $g \circ f$.
\end{proposition}

\begin{proof}
  Write $X = \bigoplus_{\ell \in L} K(u_\ell, v_\ell)$, $Y = \bigoplus_{j \in J} K(c_j, d_j)$, and $Z = \bigoplus_{i \in I} K(a_i, b_i)$. Denote by $f_{j\ell}$ the restriction of $f$ to $\hom(K(u_\ell, v_\ell), K(c_j, d_j))$ via the isomorphism \eqref{C!E:iso1}, and similarly for $g_{ij}$ and $(g \circ f)_{i\ell}$. Note the swapped indices, which are chosen to match matrix notations. Then, for all $(i, \ell) \in I \times L$, we have the equality of morphisms
  \[
    (g \circ f)_{i \ell} = \sum_{j \in J} g_{ij} \circ f_{j \ell}.
  \]
  Denote by $H \in K^{I \times L}$ the matrix representation of $g \circ f$ with columns in the same order as $F$ and rows in the same order as $G$. From the previous equality, the isomorphism \eqref{C!E:iso2}, and the definition of overlap multiplication, we deduce that
  \[
    H_{i \ell} = \sum_{j \in J} G_{ij} F_{j \ell} \bb{I}\{a_i \leq u_\ell < b_i \leq v_\ell \} = (GH)_{i \ell},
  \]
  and this concludes the proof.
\end{proof}

\subsection{Matchings}

Set-theoretic \textbf{matchings} are defined as partial injective functions between sets \cite{BauerLesnick2020}, but we will use the following equivalent definition in order to more easily generalize to multiset matchings.

\begin{definition}[Matchings]
  Let $I$ and $J$ be sets. A \index{matching, partial}(\textbf{set-theoretic partial}) \textbf{matching between $I$ and $J$} is a subset $\mu \subseteq I \times J$ such that the projections $\pi_I$ and $\pi_J$ onto $I$ and $J$, respectively, are injective. We denote by \index{.b match@$\Match$}$\Match(I, J)$ the set of matchings between $I$ and $J$.
\end{definition}

We use the definitions of the category $\mathbf{Mch}$ (including composition of matchings $\circ$), \textbf{multiset representations}, \textbf{barcodes}, and the category $\Barc$ (including the \index{composition!of barcode matchings}overlap composition of matchings $\bullet$) from \cite{BauerLesnick2020}. We also define multiset matchings:

\begin{definition}[Multiset matchings {\cite[Definition 2.4]{GST2025}}]
  Let $(S, m_S)$ and $(T, m_T)$ be two multisets. A \index{matching, partial!for multisets}(\textbf{partial multiset}) \textbf{matching between $S$ and $T$} is a multiset $(M, m_M)$ where $M = S \times T$ and, for all $s \in S$ and $t \in T$,
  \[
    \sum_{t' \in T} m_M(s, t') \leq m_S(s)
    \quad \text{and} \quad
    \sum_{s' \in S} m_M(s', t) \leq m_T(t)
    .
  \]
  We denote by $\Match(S, T)$ the set of matchings between $S$ and $T$.
\end{definition}

In particular, this definition coincides with the set-theoretic one when $S$ and $T$ are sets ($m_S$ and $m_T$ bounded above by $1$).

\begin{definition}[Matching costs]
  Let $S$ and $T$ be multisets of intervals of the real line, $M \in \Match(S, T)$ a matching, and $p \in [1, \infty]$. The \index{cost (of a matching)}\textbf{$p$-cost of $M$} is
  \[
    c_p(M) \coloneq \left( \sum_{(\langle s^\sigma, t^\tau \rangle, \langle s^{\prime \sigma'}, t^{\prime \tau'} \rangle) \in M} m_M(\langle s^\sigma, t^\tau \rangle, \langle s^{\prime \sigma'}, t^{\prime \tau'} \rangle) (|s - s'|^p + |t - t'|^p) \right)^{\frac{1}{p}}.
  \]
  In other words, $c_p(M)$ is the $p$-norm of the vector of lengths of intervals in the symmetric differences of matched intervals in $M$.
\end{definition}

\subsection{Dualization of persistence modules and matchings}
\label{C!SS:dual}

In this paper, we will mainly prove results for monomorphisms in $\Pfd$. Analogous results hold for epimorphisms by a duality argument. Specifically, we consider the contravariant dualization endofunctor $(\cdot)^* \colon \Pfd \to \Pfd$ defined by \cite[Section 3.4]{BauerLesnick2015}, which send a map $f \colon X \to Y$ to the \index{dualization}\textbf{dual map} $f^* \colon Y^* \to X^*$. In particular, this functor it sends epimorphisms to monomorphisms, and vice versa. It also sends the bar $K(a, b)$ to the bar $K(-b, -a)$, where the negative sign acts on the decorated real line $\bb{E}$ by $-(t^\tau) = (-t)^\tau$.

We also define the dual of a matching as follows.

\begin{definition}
Let $X$ and $Y$ be persistence modules in $\Pfd$ and $\mu$ a matching in $\Match(\cc{B}(X), \cc{B}(Y))$. The \textbf{dual of $\mu$} is defined as the matching $\mu^*$
\[
  \mu^* \coloneq \{ (\langle -d, -c \rangle, \langle -b, -a \rangle) \mid (\langle a, b \rangle, \langle c, d \rangle) \in \mu \}
\]
in $\Match(\cc{B}(Y^*), \cc{B}(X^*))$.
\end{definition}

\section{Bar-to-bar morphisms and induced matchings}
\label{C!S:btb}

\subsection{Bar-to-bar morphisms}

\begin{definition}
\label{C!D:btb}
  A morphism $f \colon X \to Y$ of persistence modules in $\Bar$ is \index{bar-to-bar morphism!strictly}\textbf{strictly bar-to-bar}
  if, writing $X = \bigoplus_{i \in I} K(a_i, b_i)$ and $Y = \bigoplus_{j \in J} K(c_j, d_j)$, there exists a (set-theoretic) matching $\mu \in \Match(I, J)$ such that
  \begin{equation}
  \label{C!eq:btb}
    f = \bigoplus_{(i, j) \in \mu} f_{i, j} \oplus \bigoplus_{i \in I \setminus \pi_I(\mu)} g_i \oplus \bigoplus_{j \in J \setminus \pi_J(\mu)} h_j , 
  \end{equation}
  where each $f_{ij}$ is a nonzero morphism $K(a_i, b_i) \to K(c_j, d_j)$, each $g_i$ denotes the zero morphism $K(a_i, b_i) \to 0$ and each $h_j$ denotes the zero morphism $0 \to K(c_j, d_j)$. The matrix associated to a strictly bar-to-bar morphism has at most one non zero entry for each column and one non zero entry for each row.

  A morphism $f \colon X \to Y$ in $\Pfd$ is \index{bar-to-bar morphism}\textbf{bar-to-bar} if there exist barcode decompositions of $X$ and $Y$ such that the associated morphism in $\Bar$ is strictly bar-to-bar. In other words, there exist barcode decompositions $(\varphi_X,\bigoplus_i K(a_i,b_i))$, $(\varphi_Y,\bigoplus_j K(c_j,d_j))$ of $X$ and $Y$ respectively and $g \colon \bigoplus_i K(a_i,b_i) \rightarrow \bigoplus_j K(c_j,d_j)$ strictly bar-to-bar such that the following diagram commutes:
  \[
    \begin{tikzcd}[baseline=default]
      \bigoplus_i K(a_i, b_i) \arrow[d, "g"']
      & X \arrow[d, "f"]  \arrow[l, "\varphi_X"'] \\
      \bigoplus_j K(c_j, d_j)
      & Y  \arrow[l, "\varphi_Y"]
    \end{tikzcd}
    \qedhere
  \]
\end{definition}

Not every morphism is bar-to-bar: see \cref{C!fig:non_bar_bar_morphism} for an example where no isomorphisms would make $f \colon X \to Y$ strictly bar-to-bar.

\begin{figure}[tbp]
  \centering
  \begin{subfigure}{0.6\textwidth}
    \centering
    \begin{tikzpicture}[scale = 0.4]
      \draw[line width=1pt]
        (2, 4) node[left]{2} -- (8, 4) node[right]{4}
        (0, 1) node[left]{0} -- (8, 1) node[right]{4}
        (1, 0) node[left]{1} -- (7, 0) node[right]{3}
      ;
      
      \draw[->] (3, 3.8) -- (3, 1.2);
      \draw[->] (5, 3.8) -- (5, .2);
      
      \node (X) at (10, 3.5) {$X$};
      \node (Y) at (10, .5) {$Y$};
      \draw[->] (X) -- node[right] {$f$} (Y);
    \end{tikzpicture}
  \end{subfigure}

  \vspace{1em}

  \begin{subfigure}{0.45\textwidth}
    \centering
    \begin{tikzpicture}[scale = 0.4]
      \draw[line width=1pt]
        (2, 4) node[left]{2} -- (8, 4) node[right]{4}
        (2, 0) node[left]{2} -- (8, 0) node[right]{4}
      ;
      
      \draw[->] (5, 3.8) -- (5, 0.2);
      
      \node (X) at (10, 3.5) {$X$};
      \node (Y) at (10, .5) {$X$};
      \draw[->] (X) -- node[right] {$\varphi_X$} (Y);
    \end{tikzpicture}
  \end{subfigure}
  \hfill
  \begin{subfigure}{0.45\textwidth}
    \centering
    \begin{tikzpicture}[scale = 0.4]
      \draw[line width=1pt]
      
        (0, 4) node[left]{0} -- (8, 4) node[right]{4}
        (1, 3) node[left]{1} -- (7, 3) node[right]{3}
        (0, 1) node[left]{0} -- (8, 1) node[right]{4}
        (1, 0) node[left]{1} -- (7, 0) node[right]{3}
      ;
      
      \draw[->] (3, 3.8) -- (3, 1.2);
      \draw[->] (6, 2.8) -- (6,0.2);
      
      \node (X) at (10, 3.5) {$Y$};
      \node (Y) at (10, .5) {$Y$};
      \draw[->] (X) -- node[right] {$\varphi_Y$} (Y);
    \end{tikzpicture}
  \end{subfigure}

  \caption{\textbf{Top:} example of a non bar-to-bar morphism $f : X \rightarrow Y$. \textbf{Bottom:} the only valid automorphisms $\varphi_X \colon X \rightarrow X$ and $\varphi_Y \colon Y \rightarrow Y$. In both cases, the image of the restriction of the automorphism to a bar is contained within that same bar. Consequently, they do not alter the morphism when composed.}
  \label{C!fig:non_bar_bar_morphism}
\end{figure}

\begin{proposition}
  The composition of two strictly bar-to-bar monomorphisms is strictly bar-to-bar.
\end{proposition}

\begin{proof}
  Consider strictly bar-to-bar morphisms $f \colon X \to Y$ and $f' \colon Y \to Z$ with decompositions
  \begin{align*}
    f & = \bigoplus_{(i, j) \in \mu} f_{ij} \oplus \bigoplus_{i \in I \setminus \pi_I(\mu)} g_i \oplus \bigoplus_{j \in J \setminus \pi_J(\mu)} h_j \\
    \text{and} \quad
    f' & = \bigoplus_{(j, \ell) \in \mu'} f'_{j\ell} \oplus \bigoplus_{j \in J \setminus \pi_J(\mu')} g'_j \oplus \bigoplus_{\ell \in L \setminus \pi_L(\mu)} h'_\ell,
  \end{align*}
  where in particular $\mu \in \Match(I, J)$ and $\mu' \in \Match(J, L)$ are set-theoretic matchings. Their composition $f' \circ f$ has the decomposition
  \[
    f' \circ f = \bigoplus_{\substack{(i, \ell) \in \mu' \circ \mu \\ (i, j) \in \mu, \; (j, \ell) \in \mu'}} f'_{j\ell} \circ f_{ij} \oplus \bigoplus_{i \in I \setminus \pi_I(\mu)} g_i \oplus \bigoplus_{\ell \in L \setminus \pi_L(\mu')} h'_{\ell},
  \]
  where $\mu' \circ \mu$ is the set-theoretic matching that results from composing $\mu$ and $\mu'$ as multiset matchings, and the $j$ in the first sum is unique, given $(i, \ell) \in \mu' \circ \mu$. We observe that $f'_{j\ell} \circ f_{ij}$ is nonzero if, and only if, $(i, \ell)$ is in the overlap composition $\mu' \bullet \mu$ (see \cref{C!R:overlap}). If not, then $i \in I \setminus \pi_I(\mu' \bullet \mu)$, $\ell \in L \setminus \pi_L(\mu' \bullet \mu)$, and $f'_{j\ell} \circ f_{ij}$ is equal to the sum of zero maps $(K(a_i, b_i) \to 0) \oplus (0 \to K(e_\ell, f_\ell))$. Thus we can write
  \[
    f' \circ f = \bigoplus_{\substack{(i, \ell) \in \mu' \bullet \mu \\ (i, j) \in \mu, \; (j, \ell) \in \mu'}} f'_{j\ell} \circ f_{ij} \oplus \bigoplus_{i \in I \setminus \pi_I(\mu' \bullet \mu)} g_i \oplus \bigoplus_{\ell \in L \setminus \pi_L(\mu' \bullet \mu)} h'_{\ell},
  \]
  which is an equation of the form \eqref{C!eq:btb}, and this concludes the proof.
\end{proof}


\begin{remark}
\label{C!R:no-comp}
  The composition of arbitrary bar-to-bar monomorphisms is not necessarily bar-to-bar. For example, consider the following composition of morphisms (see \cref{C!F:no-comp}):
  \[
    K(1, 4) \oplus K(2, 3) \xrightarrow{\begin{bmatrix} 1 & 0 \\ 0 & 1 \end{bmatrix}} K(1, 4) \oplus K(1, 3)
    \xrightarrow{\begin{bmatrix} 1 & 0 \\ 1 & 1 \end{bmatrix}} K(0, 4) \oplus K(1, 3),
  \]
  where the sign decorations of the start- and end-points are not important.
  Clearly the left map is bar-to-bar. The right map is bar-to-bar by doing a base change on the middle persistence module. However, the composition is a morphism between persistence modules with trivial automorphism groups, but its matrix representation looks like that of the right morphism, and therefore the composition morphism is not bar-to-bar.
\end{remark}

\begin{figure}[tbp]
  \centering
  \begin{tikzpicture}[scale = 0.4]
    \draw[line width=1pt]
      (2, 7) node[left]{$1$} -- (8, 7) node[right]{$4$}
      (4, 6) node[left]{$2$} -- (6, 6) node[right]{$3$}
      (2, 4) node[left]{$1$} -- (8, 4) node[right]{$4$}
      (2, 3) node[left]{$1$} -- (6, 3) node[right]{$3$}
      (0, 1) node[left]{$0$} -- (8, 1) node[right]{$4$}
      (2, 0) node[left]{$1$} -- (6, 0) node[right]{$3$}
    ;

    \draw[->] (7.6, 6.8) -- (7.6, 4.2);
    \draw[->] (5.6, 5.8) -- (5.6, 3.2);
    \draw[->] (7.6, 3.8) -- (7.6, 1.2);
    \draw[->] (5.6, 2.8) -- (5.6, .2);
    \draw[->] (4, 3.8) -- (4, .2);
  \end{tikzpicture}
  \caption{Representation of the morphisms in $\Bar$ of \cref{C!R:no-comp}.}
  \label{C!F:no-comp}
\end{figure}

We can view the matching $\mu$ of \cref{C!D:btb} as a matching of barcodes. Specifically, $\mu$ induces the multiset matching $(M, m_M)$ where $M = \cc{B}(X) \times \cc{B}(Y)$ and, for all $(\langle a, b \rangle, \langle c, d \rangle) \in M$,
\[
  m_M(\langle a, b \rangle, \langle c, d \rangle) \coloneq \#\{(i, j) \in \mu \mid (a_i, b_i) = (a, b) \text{ and } (c_j, d_j) = (c, d)\}.
\]
It turns out that this matching is uniquely determined as soon as we have a bar-to-bar morphism in $\Pfd$:

\begin{proposition}
\label{C!P:unique-induced-matching}
  Let $f \colon X \to Y$ be a bar-to-bar morphism in $\Pfd$, and let:
  \begin{itemize}
    \item $\varphi_X \colon X \cong \bigoplus_{i \in I} K(a_i, b_i)$ and $\varphi_Y \colon Y \cong \bigoplus_{j \in J} K(c_j, d_j)$ be barcode decompositions;
    \item $g \colon \bigoplus_{i \in I} K(a_i, b_i) \to \bigoplus_{j \in J} K(c_j, d_j)$ be a strictly bar-to-bar morphism in $\Bar$ isomorphic to $f$ (which exists by definition);
    \item $\mu$ be a matching in $\Match(I, J)$ associated to $g$ (see \cref{C!D:btb}).
  \end{itemize}
  Then the multiset matching $M \in \Match_0(\cc{B}(X), \cc{B}(Y))$ induced by $\mu$ does not depend on $\varphi_X$, $\varphi_Y$, $g$, nor $\mu$.
\end{proposition}

\begin{proof}
  We view $f$ as an object in the category $\Fun(\R \times \{0 < 1\}, \Vect_K)$. As such a representation, the decomposition \eqref{C!eq:btb} of $g$ becomes
  \[
    f \cong g \cong \bigoplus_{(i, j) \in \mu} \bf{R}^{\langle a_i, b_i \rangle}_{\langle c_j, d_j \rangle} \oplus \bigoplus_{i \in I \setminus \pi_I(\mu)} \bf{I}^{\langle a_i, b_i \rangle}_0 \oplus \bigoplus_{j \in J \setminus \pi_J(\mu)} \bf{I}^0_{\langle c_j, d_j \rangle},
  \]
  where the indecomposables $\bf{R}^*_*$ and $\bf{I}^*_*$ are adapted from \cite[Definition~5.2]{JNT2023} and defined as follows:
  \begin{itemize}
    \item $\bf{R}^{\langle a_i, b_i \rangle}_{\langle c_j, d_j \rangle}$, where $\langle a_i, b_i \rangle$ overlaps $\langle c_j, d_j \rangle$ above, is the representation of a nonzero morphism $K(a_i, b_i) \to K(c_j, d_j)$.
    \item $\bf{I}^{\langle a_i, b_i \rangle}_0$ is the representation of the zero morphism $K(a_i, b_i) \to 0$.
    \item Similarly, $\bf{I}^0_{\langle c_j, d_j \rangle}$ is the representation of the zero morphism $0 \to K(c_j, d_j)$.
  \end{itemize}
  By the Krull-Remak-Schimdt-Azumaya theorem \cite{Goro1950} (see also \cite{BotnanCrawleyBoevey2020}), this decomposition is unique up to permutation of the indecomposables. This implies that the matching $\mu$ is unique up to permutations of indices $i, i'$ in $I$ such that $a_i = a_{i'}$ and $b_i = b_{i'}$, and similarly for $J$. In particular, the induced matching $M$ on barcodes is unique. The fact that $M$ is an overlap matching is clear from its construction, and this concludes the proof.
\end{proof}

\subsection{The same-endpoint operator}
\label{C!sec:block_functor}

We define the same-endpoint operator on direct sums of bars, and then show that it is well-defined up to isomorphism on arbitrary persistence modules in $\Pfd$.

\begin{definition}[same-endpoint operator $\Delta$]
\label{C!D:Delta}
  We define $\Delta$ as the operator $\Bar \to \Bar$ that is the identity on objects and the projection
  \begin{align*}
    \hom(X, Y) & \cong \bigoplus_{i, j} \hom(K(a_i, b_i), K(c_j, d_j)) \\
    & \longrightarrow \bigoplus_{\substack{i, j \\ b_i \: = \: d_j}} \hom(K(a_i, b_i), K(c_j, d_j)) \\
    & \hookrightarrow \hom(X, Y)
  \end{align*}
  on morphisms, where $X = \bigoplus_i K(a_i, b_i)$ and $Y = \bigoplus_j K(c_j, d_j)$. We call $\Delta$ the \index{same-endpoint operator}\textbf{same-endpoint operator}.
\end{definition}

\begin{remark}
  The first isomorphism and the last inclusion are canonical. Thus the operator $\Delta$ is well-defined and does not depend on any arbitrary choices.
\end{remark}

\begin{example}
  Figure \ref{C!fig:same_end_operator} shows an example of applying the same-endpoint operator to a monomorphism in $\Bar$.
\end{example}

\begin{figure}[tbp]
\centering
  \begin{tikzpicture}[xscale=1.5, yscale=1, thick]
    \draw[-] (1.8,2.5) -- (2.8,2.5);
    \draw[-] (1,2.15) -- (2.1,2.15);
    \draw[-] (1.4,1.8) -- (2.8,1.8);
                
    \draw[-] (1,0.5) -- (2.1,0.5);
    \draw[-] (1.2,0.15) -- (2.8,0.15);
    \draw[-] (0.8,-0.2) -- (2.8,-0.2);
        
    \draw[->, thin] (1.2,2.05) -- (1.2,0.6);
    \draw[->, thin] (1.5,1.7) -- (1.5,0.6);
    \draw[->, thin] (1.8,1.7) -- (1.8,-0.1);
    \draw[->, thin] (2.5,2.4) -- (2.5,0.25);

    \draw[-, dashed] (2.8,2.5) -- (2.8,0.-0.2);
    \draw[-, dashed] (2.1,2.15) -- (2.1,0.5);

    \node at (0.3,1.1) {
      \begin{tikzcd}
        X \arrow[d, hook, "f" right] \\
        [0.9cm]
        Y
      \end{tikzcd}
    };

    \draw[->, thin] (3.2,1.2) -- (4.2,1.2) node[midway, above=4pt] {\( \Delta \)};

    \draw[-] (6.4,2.5) -- (7.4,2.5);
    \draw[-] (5.6,2.15) -- (6.7,2.15);
    \draw[-] (6,1.8) -- (7.4,1.8);
            
    \draw[-] (5.6,0.5) -- (6.7,0.5);
    \draw[-] (5.8,0.15) -- (7.4,0.15);
    \draw[-] (5.4,-0.2) -- (7.4,-0.2);
            
    \draw[->, thin] (5.8,2.05) -- (5.8,0.6);
    \draw[->, thin] (6.4,1.7) -- (6.4,-0.1);
    \draw[->, thin] (7.1,2.4) -- (7.1,0.25);
            
    \draw[-, dashed] (7.4,2.5) -- (7.4,-0.2);
    \draw[-, dashed] (6.7,2.15) -- (6.7,0.5);
            
    \node at (4.9,1.1) {
      \begin{tikzcd}
        X \arrow[d, hook, "\Delta f" right] \\
        [0.9cm]
        Y
      \end{tikzcd}
    };
  \end{tikzpicture}
  \caption{same-endpoint operator applied to a monomorphism $f \colon X \hookrightarrow Y$ in $\Bar$. The resulting monomorphism $\Delta f$ is obtained by setting to zero all the nonzero projections between bars with different end points.}
  \label{C!fig:same_end_operator}
\end{figure}

From the definition of $\Delta$ as a linear projection on hom-spaces, we immediately get the following results.

\begin{proposition}
\label{C!P:Delta-linear}
\label{C!P:Delta-additive}
  \begin{enumerate}
    \item The operator $\Delta$ is linear on hom-spaces. That is, for morphisms $f, g \colon X \to Y$, $\Delta(f + g) = \Delta f + \Delta g$.
    \item The operator $\Delta$ is additive. That is, for morphisms $f \colon X \to Y$ and $g \colon X' \to Y'$, we have the equality $\Delta(f \oplus g) = \Delta f \oplus \Delta g$ as morphisms $X \oplus X' \to Y \oplus Y'$.
  \end{enumerate}
\end{proposition}

\begin{proposition}
\label{C!P:Delta-functor}
  The operator $\Delta$ is a functor.
\end{proposition}

\begin{proof}
  Clearly $\Delta \id = \id$. For compatibility with composition, consider the composable morphisms $f \colon \bigoplus_i K(a_i, b_i) \to \bigoplus_j K(c_j, d_j)$ and $g \colon \bigoplus_j K(c_j, d_j) \to \bigoplus_k K(s_k, t_k)$. We write $f_{ij}$ for the corresponding morphism
  \[
    \bigoplus_i K(a_i, b_i) \xrightarrow{p_i} K(a_i, b_i) \to K(c_j, d_j) \xrightarrow{\iota_j} \bigoplus_j K(c_j, d_j)
  \]
  and similarly $g_{jk}$ and $(g \circ f)_{ik}$. Then
  \[
    g \circ f = \sum_{i, j, k} g_{jk} \circ f_{ij} = \sum_{\substack{i, j, k \\ b_i \geq d_j \geq t_k}} g_{jk} \circ f_{ij}
    \quad \text{and} \quad
    (g \circ f)_{ik} = \sum_{\substack{j \\ b_i \geq d_j \geq t_k}} g_{jk} \circ f_{ij}
  \]
  because all other summands are $0$. Moreover,
  \[
    \Delta f = \sum_{b_i = d_j} f_{ij},
    \quad
    \Delta g = \sum_{d_j = t_k} g_{jk},
    \quad \text{and} \quad
    \Delta (g \circ f) = \sum_{b_i = t_k} (g \circ f)_{ik}.
  \]
  Thus we get
  \[
    (\Delta (g \circ f))_{ik} = \sum_{\substack{j \\ b_i \geq d_j \geq t_k \\ b_i = t_k}} g_{jk} \circ f_{ij} = \sum_{\substack{j \\ b_i = d_j = t_k}} g_{jk} \circ f_{ij} = (\Delta g \circ \Delta f)_{ik},
  \]
  which shows the compatibility with composition.
\end{proof}

\begin{corollary}
\label{C!C:Delta-idempot}
  Let $f \colon X \hookrightarrow Y$ be a monomorphism in $\Bar$. Then
  \begin{enumerate}
    \item If $f$ is bar-to-bar, then it is isomorphic to $\Delta f$.
    \item If $f$ is strictly bar-to-bar, then $f = \Delta f$.
  \end{enumerate}
\end{corollary}

Saying that two morphisms $f, g \colon X \to Y$ are \textbf{isomorphic} means that there exist automorphisms $\varphi \colon X \to X$ and $\psi \colon Y \to Y$ such that $f = \psi \circ g \circ \varphi$.

\begin{proof}
  If $f$ is a strictly bar-to-bar monomorphism, then it is of the form
  \[
    \bigoplus_i (K(a_i, b_i) \hookrightarrow K(c_i, b_i)) \oplus \bigoplus_j (0 \to K(c_j, d_j))
  \]
  where $X = \bigoplus_i K(a_i, b_i)$ and $Y = \bigoplus_i K(c_i, b_i) \oplus \bigoplus_j K(c_j, d_j)$. The second statement then follows from the definition of $\Delta$.

  For the first statement, by definition, there exist automorphisms $\varphi_X \colon X \to X$ and $\varphi_Y \colon Y \to Y$ such that $\varphi_Y \circ f \circ \varphi_X^{-1}$ is strictly bar-to-bar. Then
  \[
    \varphi_Y \circ f \circ \varphi_X^{-1} = \Delta(\varphi_Y \circ f \circ \varphi_X^{-1}) = \Delta \varphi_Y \circ \Delta f \circ (\Delta \varphi_X)^{-1}.
  \]
  Since $\Delta$ is a functor, it sends isomorphisms to isomorphisms. Thus $\Delta f$ is isomorphic to $f$.
\end{proof}

\begin{remark}
  The converse of the first statement is also true, but it follows from \cref{C!T:Delta-is-btb}. The converse of the second statement is not true: see the right side of \cref{C!fig:change_basis}, where $f = \Delta f$ but $f$ is not strictly bar-to-bar.
\end{remark}

\begin{figure}[tbp]
  \centering
  \begin{subfigure}{.45\textwidth}
    \centering
    \begin{tikzpicture}[scale = 0.4]
      \draw[line width=1pt]
        (2, 4) node[left]{$x_1$} -- (8, 4)
        (4, 3) node[left]{$x_2$} -- (8, 3)
        (0, 1) node[left]{$y_1$} -- (8, 1)
        (2, 0) node[left]{$y_2$} -- (8, 0)
      ;

      \draw[->] (7, 3.8) -- (7, 1.2);
      \draw[->] (7.6, 2.8) -- (7.6, .2);

      \node (X) at (10, 3.5) {$X$};
      \node (Y) at (10, .5) {$Y$};
      \draw[->] (X) -- node[right] {$f$} (Y);
    \end{tikzpicture}
    \[
      \begin{blockarray}{c@{}c@{}*{2}c}
        && x_1 & x_2 \\
        \begin{block}{c[c@{}*{2}c]}
          y_1 && 1 & 0 \\
          y_2 && 0 & 1 \\
        \end{block}
      \end{blockarray}
    \]
  \end{subfigure}
  \begin{subfigure}{.45\textwidth}
    \centering
    \begin{tikzpicture}[scale = 0.4]
      \draw[line width=1pt]
        (2, 4) node[left]{$x_1$} -- (8, 4)
        (4, 3) node[left]{$\tilde{x}_2 \coloneq x_1 + x_2$} -- (8, 3)
        (0, 1) node[left]{$y_1$} -- (8, 1)
        (2, 0) node[left]{$y_2$} -- (8, 0)
      ;

      \draw[->] (7, 3.8) -- (7, 1.2);
      \draw[->] (7.6, 2.8) -- (7.6, .2);
      \draw[->] (6.4, 2.8) -- (6.4, 1.2);

      \node (X) at (10, 3.5) {$X$};
      \node (Y) at (10, .5) {$Y$};
      \draw[->] (X) -- node[right] {$f$} (Y);
    \end{tikzpicture}
    \[
      \begin{blockarray}{c@{}c@{}*{2}c}
        && x_1 & \tilde{x}_2 \\
        \begin{block}{c[c@{}*{2}c]}
          y_1 && 1 & 1 \\
          y_2 && 0 & 1 \\
        \end{block}
      \end{blockarray}
    \]
  \end{subfigure}
  \caption{Two representations of the same morphism $f \colon X \to Y$ in $\Pfd$, as bars with arrows and as matrices. On the left, the basis $\{x_1, x_2\}$ makes the monomorphism strictly bar-to-bar. This is not the case for the basis $\{x_1, \tilde{x}_2\}$ chosen on the right.}
  \label{C!fig:change_basis}
\end{figure}

\begin{proposition}
  Let $f \colon X \hookrightarrow Y$ be a monomorphism in $\Bar$. Then $\Delta f$ is also a monomorphism.
\end{proposition}

\begin{proof}
  Write $X = \bigoplus_i K(a_i, b_i)$ and $Y = \bigoplus_j K(c_j, d_j)$. Recall that $\Delta f$ is a direct sum of morphisms $\bigoplus_{b_i = b} K(a_i, b_i) \to \bigoplus_{d_j = b} K(c_j, d_j)$ for all $b \in \bb{E}$. Thus it suffices to prove, for all $b \in \bb{E}$, that $\Delta f$ restricted to $X^{(b)} \coloneq \bigoplus_{b_i = b} K(a_i, b_i)$ is a monomorphism.

  Fix $b \in \bb{E}$ and write $g \coloneq \Delta f|_{X^{(b)}}$. There exists $t \in \bb{R}$ such that $t^- < b$ and every bar $K(a_i, b)$ in $X$ satisfies $a_i \leq t^-$.
  If there were an infinite number of bars $K(c_j, d_j)$ in $Y$ such that $t \in \langle c_j, d_j \rangle$, then $Y_t$ would be infinite-dimensional. Since $Y$ is in $\Pfd$, this is not possible, so we conclude that there are only finitely many such bars, and therefore there exists $u \geq t$ in $\bb{R}$ such that $f$ maps bars $K(a_i, b_i)$ in $X$ with $b_i = b$ exclusively to bars $K(c_j, d_j)$ in $Y$ with $d_j = b$. Thus, at $u$, the same-endpoint operator does not affect $f|_{X^{(b)}}$: in other words, $g_u = (f|_{X^{(b)}})_u$. In particular, $g_u = (f|_{X^{(b)}})_u$ is injective.

  Next, we check that $g_s$ is injective for all $s < u$. Indeed, if $g_s$ sent a nonzero element $x \in X^{(b)}_s$ to $0 \in Y_s$, then by the transition maps we would get $X_{s \leq u}(x) \neq 0$ and $g_t(X_{s \leq u}(x)) = 0$, contradicting the injectivity of $g_u$. For $s \in \langle u^+, b \rangle$, injectivity holds because $g_s = (f|_{X^{(b)}})_s$ by the same argument as before. Thus $g$ is pointwise injective, so it is a monomorphism. This concludes the proof.
\end{proof}

\begin{proposition}
\label{C!prop:pfd-def-Df}
  Let $f \colon X \to Y$ be a morphism in $\Pfd$. Let $\varphi_X \colon X \cong \bigoplus_i K(a_i, b_i)$ and $\varphi_Y \colon Y \cong \bigoplus_j K(c_j, d_j)$ be barcode decompositions of $X$ and $Y$. Define $\Delta f \coloneq \varphi_Y^{-1} \Delta(\varphi_Y f \varphi_X^{-1}) \varphi_X$, where we identify the morphism $\varphi_Y f \varphi_X^{-1}$ in $\Pfd$ with the one in $\Bar$. Then, up to isomorphism, $\Delta f$ does not depend on the choice of $\varphi_X$ and $\varphi_Y$.
\end{proposition}

\begin{proof}
  Let $\psi_X$ and $\psi_Y$ be other choices of isomorphisms. Then
  \begin{align*}
    \psi_Y^{-1} \Delta(\psi_Y f \psi_X^{-1}) \psi_X & = \psi_Y^{-1} \Delta (\psi_Y \varphi_Y^{-1} \varphi_Y f \varphi_X^{-1} \varphi_X \psi_X^{-1}) \psi_X \\
    & \cong \Delta(\underbrace{\psi_Y \varphi_Y^{-1}}_{\in \Bar} \underbrace{\varphi_Y f \varphi_X^{-1}}_{\in \Bar} \underbrace{\varphi_X \psi_X^{-1}}_{\in \Bar}) \\
    & = \Delta(\psi_Y \varphi_Y^{-1}) \Delta(\varphi_Y f \varphi_X^{-1}) \Delta(\varphi_X \psi_X^{-1}) \tag{by functoriality in $\Bar$} \\
    & \cong \Delta(\varphi_Y f \varphi_X^{-1}) \\
    & \cong \varphi_Y^{-1} \Delta(\varphi_Y f \varphi_X^{-1}) \varphi_X.
    \qedhere
  \end{align*}
\end{proof}

\begin{remark}
  In particular, this means that we can define the same-endpoint operator $\Delta$ on $\Pfd$ up to isomorphism.
\end{remark}

\begin{theorem}
\label{C!T:Delta-is-btb}
  For all morphisms $f$ in $\Pfd$, the morphism $\Delta f$ is bar-to-bar.
\end{theorem}

\begin{remark}
  In $\Pfd$, the morphism $\Delta f$ is defined only up to isomorphism, see \cref{C!prop:pfd-def-Df}. Since the property of being bar-to-bar is an invariant of isomorphism classes of morphisms (by definition), \cref{C!T:Delta-is-btb} is well-formulated. 
\end{remark}

To show \cref{C!T:Delta-is-btb}, we will be working with infinite (unlabeled) matrices. In particular with infinite column-finite matrices (see \cref{C!D:column-finite-matrices}).

In \cite{Paraskevopoulos2012} the author presents a reduction algorithm for infinite row-finite matrices (defined analogously to column-finite ones). We proceed similarly with \cref{C!A:inf-red}: compare with the standard algorithm of persistent homology, see \cite{CSEM2006}.

\begin{algorithm}[tbp]
  \caption{Infinite column-finite (unlabeled) matrix reduction.
  $I$ denotes the infinite identity matrix, $E_{i, j}$ denotes the infinite matrix with all $0$'s except a $1$ in position $(i, j)$, and $\low(j) \coloneq \max\{i \geq 1 \mid A_{i, j} \neq 0\}$.}
  \label{C!A:inf-red}
  \KwIn{an infinite column-finite matrix $M$.}
  \KwOut{a factorization $A = U M V$ where $U$ and $V$ are upper triangular invertible and $A$ has at most one nonzero entry in each row and column.}

  $U \leftarrow I$\;
  $V \leftarrow I$\;
  $A \leftarrow M$\;
  \For{$j = 1, 2, \ldots$}{%
    \While(\tcp*[f]{reduce by columns with same low}){$\exists j' < j, \; \low(j') = \low(j)$}{%
      $\lambda \leftarrow A_{\low(j), j} / A_{\low(j'), j'}$\;
      $A_{*, j} \leftarrow A_{*, j} - \lambda A_{*, j'}$ \tcp*{equivalently, $A \leftarrow A (I - \lambda E_{j', j})$}
      $V \leftarrow V (I - \lambda E_{j', j})$\;
    }
    $i \leftarrow \low(j)$\;
    \While(\tcp*[f]{reduce nonzero entries above low}){$\exists i' < i, \; A_{i', j} \neq 0$}{%
      $\lambda \leftarrow A_{i', j} / A_{i, j}$\;
      $A_{i', *} \leftarrow A_{i', *} - \lambda A_{i, *}$ \tcp*{equivalently, $A \leftarrow (I - \lambda E_{i', i}) A$} 
      $U \leftarrow (I - \lambda E_{i', i}) U$\;
    }
  }
  \Return{$A$, $U$, $V$}
\end{algorithm}

\begin{lemma}
\label{C!L:good-alg}
  \cref{C!A:inf-red} is well-specified: it takes as input an infinite column-finite matrix $M$ and returns a factorization $A = U M V$ where $U$ and $V$ are upper triangular invertible and $A$ has at most one nonzero entry in each row and column.
\end{lemma}

\begin{proof}
  We observe that after the $j^{\th}$ step, the first $j$ columns of $V$ are fixed (they will not be modified in subsequent iterations). To show that the limit matrix $V$ is invertible, we construct its left inverse iteratively, starting with $V^{-1} = I$ and then applying $V^{-1} \leftarrow (I - \lambda E_{j', j}) V^{-1}$ every time we set $V \leftarrow V (I - \lambda E_{j', j})$. After the $j^{\th}$ step, the row transformations on $V^{-1}$ involve adding rows with index $k > j$ to other rows: this will never affect the first $j$ rows, as $V^{-1}$ is always upper triangular.

  For $U$, the argument is a little subtler. For all $j$, there exists a positive integer $J$ such that, in the fully reduced matrix $A$, for all $j' > J$, $\low(j') > \low(j)$. Thus, after the $J^{\th}$ step, the $j^{\th}$ column of $U$ is no longer affected by the row transformations that update $U$. To show that the limit matrix $U$ is invertible, we construct its right inverse iteratively, starting with $U^{-1} = I$ and then applying $U^{-1} \leftarrow U^{-1} (I - \lambda E_{i', i})$ every time we set $U \leftarrow (I - \lambda E_{i', i}) U$. After the $J^{\th}$ step, the only columns that get transformed are those with index $> \low(j)$, so in particular the $j^{\th}$ column is no longer affected by the column transformations that update $U^{-1}$.

  Finally, for $A$, we use the same arguments as for $V$ and $U$ to conclude that $A$ converges to a well-defined limit matrix.
\end{proof}

\begin{proof}[Proof of \cref{C!T:Delta-is-btb}]
  It suffices to prove the result for $f$ in $\Bar$, so we assume this. In this case, the morphism $\Delta f \colon X \to Y$ is a direct sum of morphisms between the direct summands of $X$ and $Y$ where all bars share a common end-point. $\Delta f$ is bar-to-bar if each summand is bar-to-bar, so, without loss of generality, we can assume that $X$ and $Y$ are persistence modules where all bars share a common end-point. Following \cref{C!R:ephemeral}, we can then represent $\Delta f$ as an infinite matrix $M$ with countably many columns and rows, and where the columns and rows correspond to bars and are sorted by start-point.

  We observe that every column of $M$ has a finite number of nonzero entries. Indeed, if the column corresponding to a bar $K(a, b)$ of $X$ had an infinite number of nonzero entries, then $K(a, b)$ would map to an infinite sequence $(K(c_i, b))_i$ of bars of $Y$, with $c_i \leq a$ for all $i$. Then (since we do not consider empty bars), there exists a point $t \in \langle a, b \rangle \subseteq \langle c_i, b \rangle$ for all $i$, and so $Y_t$ is infinite-dimensional, which contradicts the definition of $\Bar$.

  Thus we can apply \cref{C!A:inf-red} and \cref{C!L:good-alg} to obtain invertible upper triangular matrices $U$ and $V$ such that $U M V$ has at most one nonzero entry in each row and column: this corresponds to a strictly bar-to-bar morphism from $X$ to $Y$. Since $U$ and $V$ are upper triangular, they correspond to automorphisms of $X$ and $Y$, respectively. Therefore the morphism $f$ is isomorphic to a strictly bar-to-bar morphism. In other words, $f$ is bar-to-bar.
\end{proof}

\subsection{Induced matchings}

Recall that we use the the category $\Barc$ from \cite{BauerLesnick2020}. While for their canonically induced matchings compose naturally, we need more structure on our persistence modules to make the matchings compose in an intuitive way.

\begin{definition}[Matching operators]
\label{C!D:matching-operator}
  A \index{matching operator}\textbf{matching operator} is a function $\mu$ that takes a morphism $f \colon X \to Y$ in $\Bar$ and returns a matching (i.e., morphism) $\mu(f) \colon \cc{B}(X) \to \cc{B}(Y)$ in $\Barc$. The second argument is sometimes omitted. The operator $\mu$ is
  \begin{enumerate}
    \item \index{additive matching}\textbf{additive} if, given morphisms $f$ and $g$, $\mu(f \oplus g) = \mu(f) \sqcup \mu(g)$.
    \item \index{split matching}\textbf{split} if, for a strictly bar-to-bar morphism $f$, $\mu(f)$ is the matching of \cref{C!D:btb}\footnote{more precisely, the definition's matching is $\{(i, j) \mid ((\langle a_i, b_i \rangle, i), (\langle c_j, d_j \rangle, j)) \in \mu(f)\}$.}.
    \item \index{functorial matching}\textbf{functorial on monomorphisms} (resp.\ \textbf{epimorphisms}) if, given composable monomorphisms (resp.\ epimorphisms) $f$ and $g$, $\mu(g \circ f) = \mu(g) \bullet \mu(f)$. \qedhere
  \end{enumerate}
\end{definition}

\begin{remark}
  \begin{enumerate}
    \item Matching operators are well-defined. In particular, they are equivariant under reindexing of the barcode multiset representations \cite{BauerLesnick2020}.
    \item By \cite[Proposition 5.10]{BauerLesnick2015}, such operators are never functorial over all morphisms. Moreover, the \textbf{canonically induced matching} of that paper is functorial on monomorphisms and epimorphisms.
    \item Note that being split is equivalent to being additive and mapping morphisms $K(a, b) \to K(c, d)$ between single bars nontrivially.
    \item Using the axiom of choice, and viewing morphisms as objects in the additive category $\Fun(\R \times \{0 < 1\}, \Vect_K)$, we can define arbitrary additive matchings by choosing a matching for each indecomposable morphism. \qedhere
  \end{enumerate}
\end{remark}

We now show that being split and functorial on monomorphisms (resp.\ epipmorphisms) are mutually exclusive.

\begin{proposition}
\label{C!P:split-nand-functorial}
  A matching operator cannot be split and functorial on monomorphisms (resp.\ epimorphisms).
\end{proposition}

\begin{proof}
  We proceed by contradiction: suppose that such an operator $\mu$ exists. Let $b > a_1 > a_2 > \cdots > a_6$ be elements of $\E$, and consider the sequence
  \[
  X = K(a_1, b) \oplus K(a_2, b) \xhookrightarrow{\; f \;} Y = K(a_3, b) \oplus K(a_4, b) \xhookrightarrow{\; g \;} Z = K(a_5,b) \oplus K(a_6, b) 
  \]
  where the two monomorphisms are represented by the labeled matrices
  \[
  f = \begin{blockarray}{c cc}
  & a_2 & a_1 \\
  \begin{block}{c [cc]}
  a_4 & 1 & 1 \\
  a_3 & 1 & 0 \\
  \end{block}
  \end{blockarray},
  \quad
  g = \begin{blockarray}{c cc}
  & a_4 & a_3 \\
  \begin{block}{c [cc]}
  a_6 & 0 & 1 \\
  a_5 & 1 & 0 \\
  \end{block}
  \end{blockarray},
  \quad \text{and so} \quad
  gf = \begin{blockarray}{c cc}
  & a_2 & a_1 \\
  \begin{block}{c [cc]}
  a_6 & 1 & 0 \\
  a_5 & 1 & 1 \\
  \end{block}
  \end{blockarray},
  \]
  where the label $a_i$ corresponds to the interval $\langle a_i, b\rangle$.
  Recall that isomorphisms of morphisms correspond to certain left-to-right column transformations and bottom-to-top row transformations of their matrix representations. Adding the left column to the right column in the matrix representation of $gf$, we obtain the isomorphism
  \begin{equation}
  \label{C!E:counterex-in-proof}
    gf
    \cong
    \begin{blockarray}{c cc}
      & a_2 & a_1 \\
      \begin{block}{c [cc]}
        a_6 & 1 & 1 \\
        a_5 & 1 & 0 \\
      \end{block}
    \end{blockarray}
    =
    \underbrace{\begin{blockarray}{c cc}
      & a_4 & a_3 \\
      \begin{block}{c [cc]}
        a_6 & 1 & 0 \\
        a_5 & 0 & 1 \\
      \end{block}
    \end{blockarray}}_{\eqcolon \, h_0}
    \circ
    \begin{blockarray}{c cc}
      & a_2 & a_1 \\
      \begin{block}{c [cc]}
        a_4 & 1 & 1 \\
        a_3 & 1 & 0 \\
      \end{block}
    \end{blockarray}.
  \end{equation}
  In other words, $gf$ and $h_0f$ are isomorphic. Moreover, we can decompose $f$ as
  \[
    \begin{blockarray}{c cc}
      & a_2 & a_1 \\
      \begin{block}{c [cc]}
        a_4 & 1 & 1 \\
        a_3 & 1 & 0 \\
      \end{block}
    \end{blockarray}
    =
    \underbrace{\begin{blockarray}{c cc}
      & a_4 & a_3 \\
      \begin{block}{c [cc]}
        a_4 & 1 & 1 \\
        a_3 & 0 & 1 \\
      \end{block}
    \end{blockarray}}_{\eqcolon \, h_1}
    \circ
    \underbrace{\begin{blockarray}{c cc}
      & a_2 & a_1 \\
      \begin{block}{c [cc]}
        a_4 & 0 & 1 \\
        a_3 & 1 & 0 \\
      \end{block}
    \end{blockarray}}_{\eqcolon \, h_2}
  \]
  The morphism $h_1$ is an isomorphism, and by splitness $\mu(h_2)$ is also an isomorphism, so by functoriality (recall also that, by functoriality, $\mu$ sends isomorphisms to isomorphisms), we deduce that $\mu(f) = \mu(h_1) \circ \mu(h_2)$ is an isomorphism. Applying $\mu$ to \eqref{C!E:counterex-in-proof}, by functoriality, we get $\mu(g) \circ \mu(f) \cong \mu(h_0) \circ \mu(f)$. As we just showed that $\mu(f)$ is invertible, we deduce that $\mu(g) \cong \mu(h_0)$. But the only automorphisms of $\cc{B}(Y)$ and $\cc{B}(Z)$ are the identity matchings, so $\mu(g) = \mu(h_0)$. But this contradicts the assumption that $\mu$ is split, since $g$ and $h_0$ are different strictly bar-to-bar morphisms. We conclude that there is no induced matching representation operator that is both split and functorial on monomorphisms. The epimorphism case holds by a dual argument.
\end{proof}

We can now define our induced matching.

\begin{definition}
\label{C!D:Delta-matching}
  Given a morphism $f \colon X \to Y$ in $\Bar$, we know by \cref{C!T:Delta-is-btb} that $\Delta f$ is bar-to-bar. Specifically, \cref{C!A:inf-red} induces an overlap matching $\mu_\Delta(f) \in \Barc(\cc{B}(X), \cc{B}(Y))$.
\end{definition}

\begin{remark}
  \begin{enumerate}
    \item In the case where $X$ and $Y$ have finitely many bars, this induced matching is constructed for monomorphisms and epimorphisms in \cite[Theorems 6.11 \& 6.13]{bubenik2018algebraic}, where it is called an \textit{induced algebraic matching}.
    \item We can show that the matching operator $\mu_\Delta$ is additive. Moreover, it is split by construction. As a consequence, it is not functorial on either monomorphisms nor epimorphisms.
    \item The matching operator $\mu_\Delta$ is not very informative on arbitrary morphisms. Indeed, it is empty as soon as no bars of $X$ or $Y$ have the same end-point. \qedhere
  \end{enumerate}
\end{remark}

Instead, we may follow the strategy of \cite{BauerLesnick2015}: given a morphism $f \colon X \to Y$ in $\Bar$, consider an epi-mono factorization of $f$,
\[
  X \xrightarrow{e} \im(f) \xrightarrow{m} Y,
\]
where we define $\im(f)$ as the object in $\Bar$ that isomorphic to the image of $f$ in $\Pfd$, i.e.\ it is the sum of the bars in the barcode of the image. We would then define the \textbf{matching induced by $f$} to be the matching
\begin{equation}
\label{C!E:bad-induced-matching}
  \mu(f) \coloneq \mu_\Delta(m) \bullet \mu_\Delta(e^*)^*
\end{equation}
in $\Barc$, where $-^*$ is the dual operator for epimorphisms and matchings (see Section \ref{C!SS:dual}).

The issue with this construction is that the epi-mono factorization is only unique up to isomorphism: there exist other factorizations, which are all of the form $(\varphi^{-1} \circ e , m \circ \varphi)$ with $\varphi$ an automorphism on $\im(f)$. Indeed, we have the following example, which shows that the induced matching depends on the choice of factorization.

\begin{example}
\label{C!ex:counterex2}
  Consider the following morphism (over the field $\mathbb{F}_2$) along with an epi-mono factorization:
  \[
    \begin{tikzpicture}[scale = 0.4, baseline={([yshift=-.5ex]current bounding box.center)}]
      \draw[line width=1pt]
        (3, 7) node[left]{1} -- (9, 7) node[right]{3}
        (3, 6) node[left]{1} -- (6, 6) node[right]{2}
        (3, 1) node[left]{1} -- (6, 1) node[right]{2}
        (0, 0) node[left]{0} -- (6, 0) node[right]{2}
      ;
      
      \draw[->] (3.4, 6.8) -- (3.4, 1.2);
      \draw[->] (5.2, 5.8) -- (5.2, 1.2);
      \draw[->] (5.6, 5.8) -- (5.6, 0.2);
      
      \node (X) at (11, 6.5) {$X$};
      \node (Y) at (11, .5) {$Y$};
      \draw[->] (X) -- node[right] {$f$} (Y);
    \end{tikzpicture}
    =
    \begin{tikzpicture}[scale = 0.4, baseline={([yshift=-.5ex]current bounding box.center)}]
      \draw[line width=1pt]
        (3, 7) node[left]{$X_2$} -- (9, 7)
        (3, 6) node[left]{$X_1$} -- (6, 6)
        
        (3, 4) node[left]{$I_2$} -- (6, 4)
        (3, 3) node[left]{$I_1$} -- (6, 3)
        
        (3, 1) node[left]{$Y_2$} -- (6, 1)
        (0, 0) node[left]{$Y_1$} -- (6, 0)
      ;
      
      \draw[->] (3.4, 6.8) -- (3.4, 4.2);
      \draw[->] (5.4, 5.8) -- (5.4, 3.2);
      
      \draw[->] (3.4, 3.8) -- (3.4, 1.2);
      \draw[->] (5.2, 2.8) -- (5.2, 1.2);
      \draw[->] (5.6, 2.8) -- (5.6, 0.2);
      
      \node (X) at (11, 6.5) {$X$};
      \node (im) at (11, 3.5) {$\im(f)$};
      \node (Y) at (11, .5) {$Y$};
      \draw[->] (X) -- node[right] {$e$} (im);
      \draw[->] (im) -- node[right] {$m$} (Y);
    \end{tikzpicture}
  \]
  where, on the right, we have named all of the bars. The image $\im(f)$ consists of two bars $I_1$ and $I_2$ with identical support $\langle 1, 2 \rangle$. We have the matrix representations
  \begin{gather*}
  F = \begin{blockarray}{c cc}
    & X_1 & X_2 \\
    \begin{block}{c [cc]}
      Y_1 & 1 & 0 \\
      Y_2 & 1 & 1 \\
    \end{block}
  \end{blockarray},
  \quad E = \begin{blockarray}{c cc}
    & X_1 & X_2 \\
    \begin{block}{c [cc]}
      I_1 & 1 & 0 \\
      I_2 & 0 & 1 \\
    \end{block}
  \end{blockarray},
  \\
  \quad \text{and} \quad
  M = \begin{blockarray}{c cc}
    & I_1 & I_2 \\
    \begin{block}{c [cc]}
      Y_1 & 1 & 0 \\
      Y_2 & 1 & 1 \\
    \end{block}
  \end{blockarray}
  \cong \begin{blockarray}{c cc}
    & I_1 & I_2 \\
    \begin{block}{c [cc]}
      Y_1 & 1 & 0 \\
      Y_2 & 1 & 1 \\
    \end{block}
  \end{blockarray}
  \begin{blockarray}{c cc}
    & I_1 & I_2 \\
    \begin{block}{c [cc]}
      I_1 & 1 & 1 \\
      I_2 & 0 & 1 \\
    \end{block}
  \end{blockarray}
  = \begin{blockarray}{c cc}
    & I_1 & I_2 \\
    \begin{block}{c [cc]}
      Y_1 & 1 & 1 \\
      Y_2 & 1 & 0 \\
    \end{block}
  \end{blockarray}
  \end{gather*}
  for $f$, $e$, and $m$, respectively (recall that we are working over the field $\mathbb{F}_2$). The induced matching is
  \begin{align*}
    \mu_\Delta(m) \bullet \mu_\Delta(e^*)^* & = \{(I_1, Y_2), (I_2, Y_1)\} \bullet \{(X_1, I_1), (X_2, I_2)\} \\
    & = \{(X_1, Y_2), (X_2, Y_1)\}.
  \end{align*}
  Now consider $\varphi \colon \im(f) \to \im(f)$ the automorphism that permutes the two bars, with matrix representation
  \[
  \Phi = \begin{blockarray}{c cc}
    & I_1 & I_2 \\
    \begin{block}{c [cc]}
      I_1 & 0 & 1 \\
      I_2 & 1 & 0 \\
    \end{block}
  \end{blockarray}.
  \]
  Since $\Phi^{-1} = \Phi$, we compute
  \[
    \Phi^{-1} E = \begin{blockarray}{c cc}
      & X_1 & X_2 \\
      \begin{block}{c [cc]}
        I_1 & 0 & 1 \\
        I_2 & 1 & 0 \\
      \end{block}
    \end{blockarray}
    \quad \text{and} \quad
    M \Phi = \begin{blockarray}{c cc}
      & I_1 & I_2 \\
      \begin{block}{c [cc]}
        Y_1 & 0 & 1 \\
        Y_2 & 1 & 1 \\
      \end{block}
    \end{blockarray}
    \cong \begin{blockarray}{c cc}
      & I_1 & I_2 \\
      \begin{block}{c [cc]}
        Y_1 & 0 & 1 \\
        Y_2 & 1 & 0 \\
      \end{block}
    \end{blockarray}
  \]
  and find the induced matching
  \begin{align*}
    \mu_\Delta(m \circ \varphi) \bullet \mu_\Delta((\varphi^{-1} \circ e)^*)^*
    & = \{(I_1, Y_2), (I_2, Y_1)\} \bullet \{(X_1, I_2), (X_2, I_1)\} \\
    & = \{(X_1, Y_1), (X_2, Y_2)\} \\
    & \neq \mu_\Delta(m) \bullet \mu_\Delta(e^*)^*.
  \end{align*}
  Thus the induced matching as defined in \cref{C!E:bad-induced-matching} is not unique.
\end{example}

Induced matchings and automorphisms of persistence modules interact as follows.

\begin{proposition}
\label{C!P:up-to-perm}
  In the category $\Bar$, let $f \colon X \to Y$ be a morphism and $\varphi \colon X \to X$ and $\psi \colon Y \to Y$ automorphisms. Then there exist automorphic matchings $\sigma \colon \cc{B}(X) \to \cc{B}(X)$ and $\tau \colon \cc{B}(Y) \to \cc{B}(Y)$ such that
  \[
    \mu_\Delta(\psi \circ f \circ \varphi) = \tau \bullet \mu_\Delta(f) \bullet \sigma.
  \]
\end{proposition}

\begin{proof}
  The morphisms $f$ and $\psi \circ f \circ \varphi$ are isomorphic to each other. By \cref{C!prop:pfd-def-Df}, $\Delta(f)$ and $\Delta(\psi \circ f \circ \varphi)$ are isomorphic. By definition of $\mu_\Delta$ and the arguments of the proof of \cref{C!P:unique-induced-matching} (the induced matching is read from the decomposition of $f$ as a representation of $\R \times \{0 < 1\}$), the induced matchings $\mu_\Delta(f)$ and $\mu_\Delta(\psi \circ f \circ \varphi)$ are equal up to permutations of bars with identical support. By the overlap condition for composition in $\Barc$, automorphic matchings are exactly the bijective matchings that, at most, only permute of bars with identical support. We conclude that there exist automorphic matchings $\sigma \colon \cc{B}(X) \to \cc{B}(X)$ and $\tau \colon \cc{B}(Y) \to \cc{B}(Y)$ such that
  \[
    \mu_\Delta(\psi \circ f \circ \varphi) = \tau \bullet \mu_\Delta(f) \bullet \sigma.
    \qedhere
  \]
\end{proof}

The automorphic matchings $\sigma$ and $\tau$ depend on the morphism $f$, as shown by the following example.

\begin{example}
  Let $f$ and $g \colon X \to Y$ be morphisms in $\Bar$ represented by the matrices
  \[
    F = \begin{bmatrix}
      1 & 0 \\
      1 & 1
    \end{bmatrix}
    \cong \begin{bmatrix}
      1 & 1 \\
      1 & 0
    \end{bmatrix}
    \quad \text{and} \quad
    G = \begin{bmatrix}
      0 & 1 \\
      1 & 0
    \end{bmatrix},
  \]
  respectively, where the two bars of $X$ have the same support, and the two bars of $Y$ have the same support (but not necessarily the same as that of $X$). Let $\varphi$ be the automorphism on $X$ represented by
  \[
    \Phi = \begin{bmatrix}
      0 & 1 \\
      1 & 0
    \end{bmatrix}.
  \]
  Then
  \[
    F \Phi =\begin{bmatrix}
      0 & 1 \\
      1 & 1
    \end{bmatrix}
    \cong \begin{bmatrix}
      0 & 1 \\
      1 & 0
    \end{bmatrix}
    \quad \text{and} \quad
    G \Phi = \begin{bmatrix}
      1 & 0 \\
      0 & 1
    \end{bmatrix}.
  \]
  Reading the pivots of the relevant matrices, we find that $\mu_\Delta(f) = \mu_\Delta(f \varphi)$ but $\mu_\Delta(g) \neq \mu_\Delta(g \varphi)$. In other words, the automorphic matching $\sigma$ induced by $\varphi$ via \cref{C!P:up-to-perm} depends on the morphism, with which it is composed.
\end{example}

\subsection{Example: recovering the \texorpdfstring{$\mathcal{P}$}{P}-matching}
\label{C!SS:P-matching}
\index{.c pmatching@$\cc{P}$-matching}

We can define an induced matching on $\Pfd$ as soon as we determine a strategy for selecting, for a given morphism $f$ in $\Pfd$, an epi-mono factorization $f = me$. As an example, in this section we recover the induced matchings of \cite{gonzalez2020additive,GST2025} via epi-mono factorizations.

\textbf{In this section we assume that our persistence modules have countably many bars.} For example, this happens if we assume that they have no ephemeral bars.

We provide a simpler method of computing the image of a morphism in $\Bar$; compare with \cite{TorrasCasas2023}.

\begin{proposition}
\label{C!P:image-barcode}
  Let $f \colon X \to Y$ be a morphism in $\Bar$ and write $X = \bigoplus_{j \in J} K(a_j, b_j)$ and $Y = \bigoplus_{i \in I} K(c_i, d_i)$. We define the strictly bar-to-bar epimorphism
  \[
    \tilde{X} \coloneq \bigoplus_{j \in J} K(a_j, \infty) \twoheadrightarrow X
  \]
  and the strictly bar-to-bar monomorphism
  \[
    Y \hookrightarrow \tilde{Y} \coloneq \bigoplus_{i \in I} K(-\infty, d_i).
  \]
  This gives us the morphism
  \[
    \tilde{f} \coloneq \tilde{X} \twoheadrightarrow X \xrightarrow{f} Y \hookrightarrow \tilde{Y}.
  \]
  Let $\tilde{F}$ be an ordered matrix representation of $\tilde{f}$.
  Then we can apply \cref{C!A:labeled-red} to reduce the matrix $\tilde{F}$ to a matrix $A$ with at most one nonzero coefficient in each row and column. Moreover, these pivots of $A$ give the barcode of $\im(f)$:
  \[
    \cc{B}(\im(f)) = \{ \langle a_j, d_i \rangle \mid A_{ij} \neq 0 \}.
  \]
\end{proposition}

\begin{algorithm}[tbp]
  \caption{Infinite column-finite labeled matrix reduction. Notations are as in \cref{C!A:inf-red}.}
  \label{C!A:labeled-red}
  \KwIn{an infinite column-finite labeled matrix $M$ where the sequence of row labels and the sequence of column labels are totally ordered.}
  \KwOut{a factorization $A = U M V$ where $U$ and $V$ are upper triangular invertible and $A$ has at most one nonzero entry in each row and column.}

  $U \leftarrow I$\;
  $V \leftarrow I$\;
  $A \leftarrow M$\;
  \For{$j = 1, 2, \ldots$}{%
    \While(\tcp*[f]{reduce by columns with same low}){$\exists j' < j, \; \low(j') = \low(j)$}{%
      $\lambda \leftarrow A_{\low(j), j} / A_{\low(j'), j'}$\;
      $A \leftarrow A (I - \lambda E_{j', j})$\;
      $V \leftarrow V (I - \lambda E_{j', j})$\;
    }
    $i \leftarrow \low(j)$\;
    \While(\tcp*[f]{reduce nonzero entries above low}){$\exists i' < i, \; A_{i', j} \neq 0$}{%
      $\lambda \leftarrow A_{i', j} / A_{i, j}$\;
      $A \leftarrow (I - \lambda E_{i', i}) A$\;
      $U \leftarrow (I - \lambda E_{i', i}) U$\;
    }
  }
  \Return{$A$, $U$, $V$}
\end{algorithm}

\begin{proof}
  First, we observe that $\im(f) \cong \im(\tilde{f})$. Next, we recall that we have assumed that $X$ and $Y$ have countably many bars. Moreover, the row labels (resp.\ column labels) of $\tilde{F}$ are totally ordered, as they all share the same start-point (resp.\ end-point). We can therefore apply \cref{C!A:labeled-red} with $\tilde{F}$ as input. By the assumption that the row and column labels are totally ordered, we know that all left-to-right column operations and bottom-to-top row operations on the matrix $\tilde{F}$ are valid matrix operations.

  Note that \cref{C!A:labeled-red} is very similar to \cref{C!A:inf-red}, the difference being that we are reducing a labeled matrix instead of an unlabeled matrix. In particular, we use matrix multiplications, specifically overlap multiplications (\cref{C!D:labeled-matrix}), to denote the column and row operations. The proof of correctness of the algorithm is analogous to \cref{C!L:good-alg}. Thus we can apply \cref{C!A:labeled-red} to obtain a matrix $A$. As in \cref{C!L:good-alg}, this matrix has at most one nonzero coefficient in each row and column. In particular, we obtain the barcode of the image of the morphism that the matrix $A$ represents: explicitly, this barcode is
  \[
    \{ \langle a_j, d_i \rangle \mid A_{ij} \neq 0 \}.
  \]
  Moreover, since $A$ is a reduction of $\tilde{F}$, the image of $A$ is isomorphic to $\im(\tilde{f}) \cong \im(f)$. This concludes the proof.
\end{proof}

\begin{remark}
  The morphism $\tilde{f} \colon \tilde{X} \to \tilde{Y}$ is a \index{fringe presentation}\textbf{fringe presentation} of $\im(f)$ as defined in \cite{Miller2025}. We also note that the ordered matrix representation of $\tilde{F}$ is exactly that of $F$, but with relabeled rows and columns. Thus \cref{C!P:image-barcode} provides an effective way of computing the barcode of the image of a morphism, provided a matrix representation of the morphism.
\end{remark}

Let $f \colon X \to Y$ be a morphism in $\Bar$. Denote by $F$ the matrix representation of $f$ where rows are ordered (lexicographically) by end-point then start-point and the columns by start-point then end-point. Applying \cref{C!P:image-barcode}, we obtain a matrix $\tilde{F}$ with the same coefficients as $F$ but different labels, and $A = U \tilde{F} V$ its reduced form, with $U$ and $V$ invertible upper triangular matrices.

Let $R \coloneq U^{-1} A$ and $\tilde{I}$ be the subset of indices in $I$ such that row $i$ of $A$ (and $R$) contains a pivot.
Let $E$ be a matrix whose columns are labeled like $F$ and whose rows are indexed by $I$. We label row $i$ of $E$ by the interval $\langle a_j, d_i \rangle$, corresponding to the pivot $(i, j)$ of $A$. We set
\[
  E_{ij} \coloneq \begin{cases}
    R_{ij} & \text{if } i = \low(j) \text{ in } R, \\
    0 & \text{otherwise}.
  \end{cases}
\]

\begin{proposition}
  We use the notations introduced immediately above. Let $Z$ be the persistence module in $\Bar$ that is isomorphic to $\im(f)$.
  There exists an epi-mono factorization $f = m \circ e$ where the epimorphism $e \colon X \to Z$ is represented by the matrix $E V^{-1}$ and the monomorphism $m \colon Z \to Y$ is represented by a matrix $M$ in row echelon form with pivots $\{(i, i) \mid i \in \tilde{I}\}$.
\end{proposition}

\begin{proof}
  We see that $E V^{-1}$ represents an epimorphism from $X$ to $\im(f)$.

  Let us solve the equation $M E = R \coloneq U^{-1} A$ for the matrix $M$. Note that the labeling of $M$ is determined by this equation. By construction, $E$ has at most one nonzero coefficient in each column. Therefore, a given unknown $M_{i, i'}$ in $M$ appears in at most one equation:
  \begin{itemize}
    \item If there exists $j \in J$ such that $i' = \low(j)$ in $R$, then $E_{i', j}$ is nonzero and $M_{i, i'} = M_{i, i'}$ must satisfy the equation
    \[
      M_{i, i'} E_{i', j} = R_{ij}.
    \]
    If $R_{ij} = 0$, then we set $M_{i, i'} = 0$. Otherwise, we need to check that we can set $M_{i, i'}$ to be nonzero. We know that $i \leq i'$ ($R_{i', j}$ is the pivot, so $R_{ij}$ must be above it). Column $i'$ of $M$ is labeled by $\langle a_j, d_{i'} \rangle$ and row $i$ of $M$ is labeled by $\langle c_i, d_i \rangle$. We have $d_{i'} \geq d_i$ because $i' \geq i$ and we assume that the $\langle c_i, d_i \rangle$ are ordered compatibly with the product order. We also have $a_j \geq c_i$ because $R_{ij}$ is nonzero (and so $\langle a_j, b_j \rangle$ must overlap $\langle c_i, d_i \rangle$ above). Thus we have checked that $\langle a_j, d_{i'} \rangle$ overlaps $\langle c_i, d_i \rangle$ above, so we can set $M_{i, i'}$ to be nonzero.
    \item Otherwise, $M_{i, i'}$ has no equations to satisfy, and we choose to set $M_{i, i'} = 0$.
  \end{itemize}
  We observe that the resulting $M$ is in reduced row echelon form, with pivots at the coordinates $(i, i)$ for $i$ in $\tilde{I}$. Each column as a pivot and the pivots' row and column labels have the same end-point $d_i$, so we deduce that $M$ represents a monomorphism: indeed, at every time $t \in \R$, $M_t$ represents an injective linear map.

  We thus obtain $F = M E V^{-1}$, giving us an epi-mono factorization $f = m e$ by the morphisms $e$ and $m$ represented by $E V^{-1}$ and $M$, respectively.
\end{proof}

\begin{definition}
  We define the induced matching $\mu(f) \coloneq \mu_\Delta(m) \bullet \mu_\Delta(e^*)^*$.
\end{definition}

\begin{remark}
\label{C!R:knowing-epi-suffices}
  By the last observation on $M$, we deduce that the induced matching $\mu_\Delta(m)$ matches each bar $\langle a_j, d_i \rangle$ in the image to $\langle c_i, d_i \rangle$ in $\cc{B}(Y)$. In other words, we know the matching $\mu(f)$ as soon as we know $\mu_\Delta(e^*)^*$, or as soon as we know the reduced matrix $A$.
\end{remark}

In the case where $X$ and $Y$ are tame persistence modules (i.e., they have finite barcodes), \cite[Algorithm 1]{GST2025} is a matrix reduction algorithm to obtain a multiset matching from a morphism $f \colon X \to Y$, which we will call the $\cc{P}$-matching.

We first show the pairing uniqueness lemma \cite{CSEM2006} for tame fringe presentations, following the reasoning and notations of the original proof. Given a matrix $A \in K^{m \times n}$, for all $i \in \{1, \ldots, m\}$ and $j \in \{1, \ldots, n\}$, denote by $A_i^j$ the lower left minor of $A$ obtained by removing the first $i - 1$ rows and the last $n - j$ columns. Define
\[
  r_A(i, j) \coloneq \rank A_i^j - \rank A_{i + 1}^j - \rank A_i^{j - 1} + \rank A_{i + 1}^{j - 1},
\]
where $\rank A$ denotes the number of bars in the image of the morphism represented by $A$.

\begin{lemma}[Pairing uniqueness lemma]
\label{C!L:pairing-uniqueness}
  Let $\tilde{F} \in K^{m \times n}$ be the relabeled matrix representation of a morphism $f \colon X \to Y$ as in \cref{C!P:image-barcode}. Now consider the reduced matrix $A$ returned by \cref{C!A:labeled-red}. Then $A_{i, j} \neq 0$ if, and only if, $r_{\tilde{F}}(i, j) = 1$.
\end{lemma}

\begin{proof}
  Note that left-to-right column operations and bottom-to-top row operations correspond to isomorphism actions on the domain and codomain of the represented morphism. This holds for the restricted matrices $A_i^j$, and therefore $r_{\tilde{F}} = r_A$. It therefore suffices to show that $A_{i, j} \neq 0$ if, and only if, $r_A(i, j) = 1$, and this follows immediately from the fact that $A$ has at most one nonzero coefficient in each row and column.
\end{proof}

\begin{proposition}
  Let $\tilde{F}$ be the relabeled matrix representation of a morphism $f \colon X \to Y$ as in \cref{C!P:image-barcode}. Then the reduction of the matrix $\tilde{F}$ by \cref{C!A:labeled-red} and the reduction of the matrix $F$ by \cite[Algorithm 1]{GST2025} have the same pivots.
\end{proposition}

\begin{proof}
  By \cref{C!L:pairing-uniqueness}, the order in which we do column and row operations in \cref{C!A:labeled-red} does not affect the resulting pivots. Moreover, we note that the row operations do not affect the pivots of the matrix, and therefore we will ignore that part of \cref{C!A:labeled-red}. We can therefore assume that we do column reductions in the same order as in \cite[Algorithm 1]{GST2025}: that is, start from the columns with the lowest nonzero coefficient and reduce them before moving to the columns with the next lowest nonzero coefficient, and so on.

  The only difference between the two algorithms is then how invalid matrix coefficients are treated. In \cref{C!A:labeled-red}, a coefficient is set to zero if the column label does not overlap above the row label; in our case where the column labels all end at $+\infty$ and the row labels start at $-\infty$, this only happens when the column and row labels are disjoint. In \cite[Algorithm 1]{GST2025}, a coefficient is set to zero if the column and row labels in $F$ are disjoint. This happens only in the following situation. There are two columns $j$ and $j'$ with $j' < j$ such that $\low(j) = \low(j') \eqcolon i$. When adding column $j'$ to column $j > j'$, a nonzero coefficient on row $i' > i$ produces an invalid nonzero coefficient at $(i', j)$. Denoting by $[a_j, b_j)$, $[a_{j'}, b_{j'})$, $[c_i, d_i)$, and $[c_{i'}, d_{i'})$ the labels of the relevant columns and rows, we have the following inequalities. By assumption on the order of the rows and columns, we have $a_{j'} \leq a_j$ and $d_{i'} < d_i$. By the validity of the nonzero coefficient at $(i', j')$, we have $c_{i'} \leq a_{j'} < d_{i'} \leq b_{j'}$. By the validity of the nonzero coefficient at $(i, j)$, we have $c_i \leq a_j < d_i \leq b_j$. Thus we have $c_{i'} < a_j$ and $d_{i'} < b_j$. In order for the coefficient at $(i', j)$ to be invalid, it is therefore necessary that $a_j < d_{i'}$ does not hold, i.e., that $a_j \geq d_{i'}$. We conclude that the same coefficient $(i', j)$ in $\tilde{F}$ would have disjoint row and column labels, and therefore would be canceled by the overlap multiplication.

  We conclude that \cref{C!A:labeled-red} and \cite[Algorithm 1]{GST2025} return matrices with the same pivots.
\end{proof}

We deduce that our induced matching $\mu(f)$ is a representation of the (multiset) $\cc{P}$-matching defined in \cite{GST2025}.

\section{Application to Wasserstein distances}
\label{C!S:Wasserstein}

\subsection{Sizes and ranks of (co)kernels}
\label{C!SS:coker}

In this section we prove some results in the category $\Pfd$. In particular, we highlight a connection between rank functions and $p$-norms, which yields new proofs of results in \cite[Section 3]{AGRS2025}. We denote by $\rk f$ the rank of a linear map $f$.

\begin{proposition}
\label{C!P:cokernel-rank}
  Let $f \colon X \hookrightarrow Y$ be a monomorphism in $\Pfd$ and consider $\Delta f$ the functorially defined bar-to-bar morphism and $f_\varphi$ the bar-to-bar monomorphism induced by the canonical matching between $X$ and $Y$ (this assumes that we fix barcodes). Then, for all $s \leq t$ in $[0, \infty)$,
  \[
    \rk((\coker f)_{s \leq t}) \geq \rk((\coker \Delta f)_{s \leq t}) \geq \rk((\coker f_\varphi)_{s \leq t}).
  \]
\end{proposition}

We first prove a series of lemmas.

\begin{lemma}
\label{C!L:coker-rank-eq}
  Let $f \colon X \hookrightarrow Y$ be an arbitrary monomorphism in $\Pfd$. For all $s \leq t$ in $[0, \infty)$,
  \[
    \rk((\coker f)_{s \leq t}) = \rk(Y_{s \leq t} + f_t) - \dim X_t,
  \]
  where $Y_{s \leq t} + f_t \colon Y_s \oplus X_t \to Y_t$.
\end{lemma}

\begin{proof}
  Consider the following diagram of vector spaces:
  \[
    \begin{tikzcd}
    X_s
      \arrow[r, hook, "f_s"]
      \arrow[d, "X_{s \leq t}" swap]
    &
    Y_s
      \arrow[r, two heads]
      \arrow[d, "Y_{s \leq t}" swap]
    &
    (\coker f)_s
      \arrow[d, "(\coker f)_{s \leq t}"]
    \\
    X_t
      \arrow[r, hook, "f_t" swap]
    &
    Y_t
      \arrow[r, two heads]
      \arrow[d, two heads]
    &
    (\coker f)_t
      \arrow[d, two heads]
    \\
    &
    \coker Y_{s \leq t}
      \arrow[r, two heads]
    &
    C
    \end{tikzcd}
  \]
  where $C$ is the cokernel of the map $(\coker f)_{s \leq t}$. We observe (by a diagram chase, for instance) that $C$ is also the cokernel of the map $g \coloneq Y_{s \leq t} + f_t \colon Y_s \oplus X_t \to Y_t$. Thus $\dim C = \dim Y_t - \rk g$, and we conclude the proof by observing that
  \begin{align*}
    \rk((\coker f)_{s \leq t}) & = \dim (\coker f)_t - \dim C \\
    & = \dim Y_t - \dim X_t - \dim Y_t + \rk g \\
    & = \rk g - \dim X_t. \qedhere
  \end{align*}
\end{proof}

\begin{lemma}
\label{C!L:delta-rank-ineq}
  Let $g \colon X \to Y$ be an arbitrary morphism in $\Pfd$. Then, for all $t$ in $[0, \infty)$,
  \[
    \rk g_t \geq \rk (\Delta g)_t.
  \]
\end{lemma}

\begin{proof}
  Fix barcode decompositions of $X$ and $Y$. Then we can represent $g_t$ as a (finite) matrix, where the columns and rows correspond to bars of $X$ and $Y$ that are nonzero at $t$, and they are sorted such that the death time is increasing. The matrix $M$ of $g_t$ is then block upper triangular, where the blocks correspond to bars with the same death time. The matrix $M_\Delta$ of $(\Delta g)_t$ is the block diagonal part of this matrix.

  Let us show that the rank of $M_\Delta$ is smaller than the rank of $M$. Indeed, for every block of $M_\Delta$, select a maximal set of linearly independent columns. The corresponding columns in $M$ are also linearly independent, and therefore the linear map sending each of the selected columns of $M_\Delta$ to the corresponding column in $M$ is injective. This defines an injective map $\im M_\Delta \hookrightarrow \im M$. Taking the dimension of both sides, we get $\rk (\Delta g)_t \leq \rk g_t$.
\end{proof}

\begin{lemma}
\label{C!L:coker-rank-ineq}
  Let $f \colon X \hookrightarrow Y$ be an arbitrary monomorphism in $\Pfd$. For all $s \leq t$ in $[0, \infty)$,
  \[
    \rk(Y_{s \leq t} + f_t) \geq \rk(Y_{s \leq t} + (\Delta f)_t),
  \]
  where $Y_{s \leq t} + f_t$ and $Y_{s \leq t} + (\Delta f)_t \colon Y_s \oplus X_t \to Y_t$.
\end{lemma}

\begin{proof}
  Since we are concerned with ranks, which are invariant by isomorphisms of morphisms, we can assume without loss of generality that $X$, $Y$, and $f$ are in $\Bar$.

  Consider $Y[-\delta]$, the persistence module defined as $Y[-\delta]_a \coloneq Y_{a - \delta}$. Then there exists a natural morphism in $\Pfd$ $Y_{\bullet - \delta \leq \bullet} \colon Y[-\delta] \to Y$.
  The image of this morphism then has the same barcode as $Y$, except that every birth time has been increased by $\delta$, with resulting bars with a birth time greater than the death time being removed. Denote by $\iota \colon \im Y_{\bullet - \delta \leq \bullet} \hookrightarrow Y$ the monomorphism through which $Y_{\bullet - \delta \leq \bullet}$ factors. Then $\rk(Y_{s \leq t} + f_t) = \rk(\iota_t + f_t)$. Moreover, $\iota$ is strictly bar-to-bar, so by \cref{C!C:Delta-idempot} $\Delta \iota = \iota$.
  By linearity of $\Delta$ (\cref{C!P:Delta-linear}), we deduce that $\Delta(\iota + f) = \iota + \Delta f$. We conclude as follows:
  \begin{align*}
    \rk(Y_{s \leq t} + f_t) & = \rk(\iota + f)_t \\
    & \geq \rk(\Delta(\iota + f))_t \tag{by \cref{C!L:delta-rank-ineq}}\\
    & = \rk(\iota_t + (\Delta f)_t) \\
    & = \rk(Y_{s \leq t} + (\Delta f)_t). \qedhere
  \end{align*}
\end{proof}

\begin{lemma}
\label{C!L:canonical-coker}
  Suppose that all of the bars of $X$ and $Y$ have the same end-point and there exists a monomorphism $X\hookrightarrow Y$ in $\Pfd$. Then
  \[
    \rk((\coker f_\varphi)_{s \leq t}) = (\rk Y_{s \leq t} - \dim X_t)_+,
  \]
  where $(-)_+$ is the positive part function.
\end{lemma}

\begin{proof}
  Denote by $b$ this common end-point. Since $f_\varphi$ is induced by the canonical matching, we can assume without loss of generality that $X = \bigoplus_{i = 1}^m K(a_i, b)$, $Y = \bigoplus_{i = 1}^n K(c_i, b)$ with $m \leq n$ in $\N \cup \{\omega\}$, the $a_i$ and $c_i$ are ordered in nondecreasing order, and the canonical matching $\varphi$ is the inclusion of $\{1, \ldots, m\}$ in $\{1, \ldots, n\}$. The fact that $X$ and $Y$ have at most countably many bars with end-point $b$ is explained in \cref{C!R:ephemeral}. Then
  \[
    f_\varphi = \bigoplus_{i = 1}^m (K(a_i, b) \to K(c_i, b)) \oplus \bigoplus_{j = m + 1}^n (0 \to K(c_i, b)).
  \]
  Therefore
  \[
    \coker f_\varphi = \bigoplus_{i = 1}^m K(c_i, a_i) \oplus \bigoplus_{j = m + 1}^n K(c_i, b).
  \]
  If $t^- \geq b$, then the result holds trivially.
  Now consider the case $t^- < b$. Since the $a_i$ and $c_i$ are in nondecreasing order, there exist integers $k_a$ and $k_c$ such that
  \begin{align*}
    \{i \in \{1, \ldots, m\} : a_i \leq t^-\} & = \{1, \ldots, k_a\} \\
    \text{and} \quad \{i \in \{1, \ldots, n\} : c_i \leq s^-\} & = \{1, \ldots, k_c\}.
  \end{align*}
  Therefore
  \begin{align*}
    & \{i \in \{1, \ldots, n\} : c_i \leq s^-\} \setminus \{i \in \{1, \ldots, m\} : a_i \leq t^-\} \\
    & = \begin{cases}
      \{k_a + 1, \ldots, k_c\} & \text{if } k_a < k_c, \\
      \varnothing & \text{otherwise},
    \end{cases}
  \end{align*}
  and we conclude that
  \begin{align*}
    \rk((\coker f_\varphi)_{s \leq t}) = {} & \big|\{i \in \{1, \ldots, m\} : c_i \leq s^-, t^- < a_i\} \\
    & \hphantom{\big|} \sqcup \{j \in \{m + 1, \ldots, n\} : c_i \leq s^-\} \big| \\
    = {} & \big| \{i \in \{1, \ldots, n\} : c_i \leq s^-\} \setminus \{i \in \{1, \ldots, m\} : a_i \leq t^-\} \big| \\
    = {} & (k_c - k_a)_+ \\
    = {} & (\rk Y_{s \leq t} - \dim X_t)_+,
  \end{align*}
  where the second equality comes from the equality between the sets, and the fourth equality comes from the fact that $k_a = \dim X_t$ and $k_c = \rk Y_{s \leq t}$.
\end{proof}

\begin{proof}[Proof of \cref{C!P:cokernel-rank}]
  We show the first inequality as follows:
  \begin{align*}
    \rk((\coker f)_{s \leq t}) & = \rk(Y_{s \leq t} + f_t) - \dim X_t \tag{by \cref{C!L:coker-rank-eq}} \\
    & \geq \rk(Y_{s \leq t} + (\Delta f)_t) - \dim X_t \tag{by \cref{C!L:coker-rank-ineq}} \\
    & = \rk((\coker \Delta f)_{s \leq t}). \tag{by \cref{C!L:coker-rank-eq}}
  \end{align*}
  We show the second inequality as follows. Since $\Delta f$ is bar-to-bar by \cref{C!T:Delta-is-btb}, its cokernel is the direct sum of the cokernels of the maps restricted to bars with a common end-point. Thus, without loss of generality, we can assume that the bars of $X$ and $Y$ all have the same end-point. In this case, we get
  \begin{align*}
    \rk((\coker \Delta f)_{s \leq t}) & = \rk(Y_{s \leq t} + (\Delta f)_t) - \dim X_t \tag{by \cref{C!L:coker-rank-eq}} \\
    & \geq \max\{\rk Y_{s \leq t}, \rk (\Delta f)_t\} - \dim X_t \tag{true for general sums of maps} \\
    & = \max\{\rk Y_{s \leq t}, \dim X_t\} - \dim X_t \tag{since $\Delta f$ is a monomorphism}\\
    & = (\rk Y_{s \leq t} - \dim X_t)_+ \\
    & = \rk((\coker f_\varphi)_{s \leq t}). \tag*{(by \cref{C!L:canonical-coker}) \qedhere}
  \end{align*}
\end{proof}

To connect rank functions to $p$-norms of cokernels, we define the following measurable space.

\begin{definition}
  Fix $p \in (1, \infty)$. We define the space $\Omega \coloneq \{(s, t) : s \leq t \in \R\} \subset \R^2$ and the density function $g_p \colon \R^2 \to [0,\infty )$ given by
  \[
    g_p(s, t) \coloneq p (p - 1) |t - s|^{p - 2}.
    \qedhere
  \]
\end{definition}

\begin{proposition}
\label{C!P:p-norm-by-integral}
  For $X$ a persistence module and $p \in (1, \infty)$,
  \[
    \int_\Omega \rk X_{s \leq t} g_p(s, t) \textnormal{d}s \textnormal{d}t = \| X \|^p_p.
  \]
\end{proposition}

\begin{proof}
  We first consider the case where $X$ is a single bar $K(a, b)$. Suppose that $a$ and $b$ are finite (i.e., not $-\infty$ nor $+\infty$, see Definition \ref{C!def:decoratedE}). Then, for all $p \in (1, \infty)$, we compute
  \begin{align*}
    \int_\Omega \rk K(a, b)_{s \leq t} g_p(s, t) \text{d}s \text{d}t & = \int_a^b \int_s^b p(p - 1)(t - s)^{p - 2} \text{d}s \text{d}t \\
    & = \int_a^b p (b - s)^{p - 1} \text{d}s \\
    & = (b - a)^p \\
    & = \| K(a, b) \|_p^p.
  \end{align*}
  If $a$ or $b$ is not finite, then we still have the formal equality $\infty = \infty$.

  If $X\cong \bigoplus_{i=1}^m K(a_i,b_i)$ is a persistence module with finitely many bars, then $\rk (X)=\sum_{i=1}^m \rk K(a_i,b_i)$ and the result follows from the case of a single bar by linearity of the Lebesgue integral. 
  Since ephemeral bars do not contribute to the integral, it remains to prove the result for persistence modules $X\cong \bigoplus_{i=1}^\infty K(a_i,b_i)$ with countably many bars (see Remark~\ref{C!R:ephemeral}).
  In this case, the equality 
  \[
    \int_\Omega \sum_{i=1}^\infty \rk K(a_i, b_i)_{s \leq t} g_p(s, t) \text{d}s \text{d}t = \sum_{i=1}^\infty \int_\Omega \rk K(a_i, b_i)_{s \leq t} g_p(s, t) \text{d}s \text{d}t
  \]
  follows from the Monotone Convergence Theorem \cite[Thm.~2.14]{folland1999real}.
\end{proof}

\begin{remark}\label{C!rmk:int-norm1}
  When $p=1$, as observed for example in \cite{BSS2023}, the $1$-norm $\| X \|_1$ of a persistence module coincides with the integral $\int_{\R} \dim X_{t} \textnormal{d}t$ of its Hilbert function $(\dim X)(t)\coloneqq \dim X_t$. 
\end{remark}

\begin{theorem}
\label{C!thm:cost-coker-btb}
  Let $f \colon X \hookrightarrow Y$ be a monomorphism in $\Pfd$ and $f_\varphi$ the bar-to-bar monomorphism induced by the canonical matching between $X$ and $Y$. Then, for all $p \in [1, \infty]$,
  \[
    \| \coker f \|_p \geq \| \coker \Delta f \|_p \geq \| \coker f_\varphi \|_p.
  \]
\end{theorem}

\begin{proof}
  For $p \in (1, \infty)$, we combine \cref{C!P:cokernel-rank,C!P:p-norm-by-integral}, along with the fact that the function $h \mapsto \int_\Omega h g_p$ is increasing.

  For $p = 1$ and $p = \infty$, we conclude by continuity of $p$-norms with respect to $p$.
\end{proof}


\begin{remark}
  When $p=1$, the inequalities of Theorem \ref{C!thm:cost-coker-btb} are equalities:
  \[
    \| \coker f \|_1 = \| \coker \Delta f \|_1 = \| \coker f_\varphi \|_1.
  \]
  To see this, recall Remark \ref{C!rmk:int-norm1} and observe that for all $t \in \R$,
  \[
    \dim (\coker f)_t = \dim (\coker \Delta f)_t = \dim (\coker f_\varphi)_t = \dim Y_t - \dim X_t.
  \]
  For all $p>1$, it is easy to find a monomorphism $f \colon X \hookrightarrow Y$ such that the inequalities of Theorem \ref{C!thm:cost-coker-btb} are strict.
\end{remark}

Dually, we obtain the following result.

\begin{theorem}
\label{C!thm:cost-ker-btb}
  Let $f \colon X \twoheadrightarrow Y$ be an epimorphism in $\Pfd$ and $f_\varphi$ the bar-to-bar epimorphism induced by the canonical matching between $X$ and $Y$. Then, for all $p \in [1, \infty]$,
  \[
    \| \ker f \|_p \geq \| \ker (\Delta f^*)^* \|_p \geq \| \ker f_\varphi \|_p.
  \]
\end{theorem}

\subsection{Recovering isometry result for Wasserstein distances}
\label{C!SS:isometry}

We now provide alternate proofs of the construction of algebraic $p$-Wasserstein distances and of the isometry result for $p$-Wasserstein distances \cite{SkrabaTurner2023}. We define algebraic $p$-Wasserstein distances by defining a noise system \cite{SCLRÖ2017}: see \cite{AGRS2025} for the details to this approach in the context of \textit{tame} persistence modules. In particular, we need to show the following result.

\begin{proposition}
  Let $0 \to X \to Y \to Z \to 0$ be a short exact sequence in $\Pfd$. Then, for all $p \in [1, \infty]$,
  \begin{enumerate}
    \item $\| X \|_p \leq \| Y \|_p$ and $\| Z \|_p \leq \| Y \|_p$ (``easy'' axiom);
    \item $\| Y \|_p \leq \| X \|_p + \| Z \|_p$ (``hard'' axiom).
  \end{enumerate}
\end{proposition}

To prove the ``easy'' axiom of noise systems we can proceed as follows, showing a stronger statement.

\begin{proposition}
\label{C!prop:rkYge}
  For every short exact sequence $0\to X\to Y\to Z\to 0$ it holds that $\rk (Y) \ge \rk (X) + \rk (Z)$.
\end{proposition}

\begin{proof}
  Let $s\le t$ in $\R$ and consider the appropriate diagram with two exact rows as in the snake lemma. We have
  \[
    \rk X_{s\le t} - \rk Y_{s\le t} + \rk X_{s\le t} = - \dim \ker X_{s\le t} + \dim \ker Y_{s\le t} - \dim \ker Z_{s\le t} \le 0,
  \]
  where the equality follows from $\rk X_{s\le t} = \dim X_t - \dim \ker X_{s\le t}$ and similar identities for $Y$ and $Z$, together with the identity $\dim X_t - \dim Y_t + \dim Z_t =0$ by exactness of $0\to X_t\to Y_t\to Z_t\to 0$; the inequality above follows from the exact sequence $0 \to \ker X_{s\le t} \to \ker Y_{s\le t} \xrightarrow{g} \ker Z_{s\le t} \to \coker g \to 0$ (easy part of the snake lemma) again by Euler characteristic.
\end{proof}

By Proposition \ref{C!P:p-norm-by-integral} and monotonicity of the integral, we have the following consequence of Proposition \ref{C!prop:rkYge}.

\begin{corollary}[{\cite[Remark 7.32]{SkrabaTurner2023}}]
  If $p \in [1, \infty)$ and $0 \to X \to Y \to Z \to 0$ is a short exact sequence, then
  \[
    \left( \| X\|_p^p + \| Z\|_p^p \right)^\frac{1}{p} \le \| Y\|_p .
  \]
\end{corollary}

The ``easy'' axiom of noise systems, which states that $\| X\|_p \leq \| Y\|_p$ and $\| Z\|_p \leq \| Y\|_p$ follows immediately. The case $p=\infty$ follows by taking the limit. 

We now prove the ``hard'' axiom of noise systems. This proof is the same as \cite[Lemma 4.15]{AGRS2025}; the novelty is proving the result for $\Pfd$.

\begin{proposition}
  For every short exact sequence $0\to X\to Y\to Z\to 0$ it holds that $\| Y \|_p \leq \| X \|_p + \| Z \|_p$.
\end{proposition}

\begin{proof}
  Denote by $f \colon X \to Y$ the monomorphism in the short exact sequence. Then we have $Z \cong \coker f$.

  First, let us consider the case where $f$ is bar-to-bar. By definition (see \cref{C!D:btb}), there exist barcode decompositions $X \cong \bigoplus_{i \in I} K(a_i, b_i)$ and $Y \cong \bigoplus_{j \in J} K_(c_j, d_j)$ and a matching $\mu \in \Match(I, J)$ such that
  \[
    \coker f \cong \bigoplus_{(i, j) \in \mu} \coker(K(a_i, b_i) \to K(c_j, d_j)) \oplus \bigoplus_{j \in J \setminus \pi_J(\mu)} K(c_j, d_j),
  \]
  where, for $(i, j)$ in $\mu$, $b_i = d_j$.
  The $p$-norm of this persistence module is unaffected by ephemeral summands, so without loss of generality we assume that there are no ephemeral summands (for countability reasons). The $p$-norm of $\coker f$ is then equal to the $p$-norm of the countably long sequence $(\ell_j)_{j \in J}$ where, for all $j \in J$,
  \[
    \ell_j = \begin{cases}
      \underbrace{(b_i - d_j)}_{= 0} {} + c_j - a_i & \text{if there exists $(i, j)$ in $\mu$}, \\
      d_j - c_j & \text{otherwise}.
    \end{cases}
  \]
  Since $f$ is a monomorphism, every $i \in I$ appears in this sequence.
  We also have $\| X \|_p = \| (b_i - a_i)_{i \in I} \|_p$ and $\| Y \|_p = \| (d_j - c_j)_{j \in J} \|_p$, and thus we obtain the desired inequality as the triangular inequality for the $p$-norm on infinite sequences of numbers.

  Now let $f$ be an arbitrary monomorphism, and let $f_\varphi$ be the bar-to-bar monomorphism induced by the canonical matching between $X$ and $Y$. By \cref{C!thm:cost-coker-btb}, we conclude with
  \[
    \| Y \|_p \leq \| X \|_p + \| \coker f_\varphi \|_p \leq \| X \|_p + \| \coker f \|_p = \| X \|_p + \| Z \|_p.
    \qedhere
  \]
\end{proof}

We provide an alternate (for one direction at least) proof to \cite[Theorems 7.27-7.28]{SkrabaTurner2023}, the isometry result between algebraic and combinatorial $p$-Wasserstein distances.

\begin{theorem}
\label{C!T:isometry}
  For any $X$ and $Y$ in $\Pfd$, we have $d_p^p (X,Y) = W_p (\mathsf{Dgm}(X),\mathsf{Dgm}(Y))$.
\end{theorem}

\begin{proof}
  \noindent\fbox{$\le$}
  Let $\varphi \colon \mathsf{Dgm}(X) \to \mathsf{Dgm}(Y)$ be a matching. We want to construct a span with a smaller or equal cost. Following \cite[Theorem 7.28]{SkrabaTurner2023}, define
  \[
    Z \coloneq \bigoplus_{((a, b), (c, d)) \in \varphi} K(\max(a, c), \max(b, d))
  \]
  and construct the natural span $X \leftarrow Z \rightarrow Y$. We can check that the cost of this span is exactly the cost of $\varphi$.

  \noindent\fbox{$\geq$}
  Let $X \xleftarrow{f} Z \xrightarrow{g} Y$ be a span. We want to construct a matching with a smaller or equal cost. First, factor the span as $X \hookleftarrow \im f \twoheadleftarrow Z \twoheadrightarrow \im g \hookrightarrow Y$ and replace each morphism with the morphism induced by canonical matchings. By \cref{C!thm:cost-coker-btb,C!thm:cost-ker-btb}, the cost of the new span $X \xleftarrow{f_b} Z \xrightarrow{g_b} Y$ is smaller than the cost of the original span.

  The (strictly bar-to-bar) morphisms $f_b$ and $g_b$ induce matchings $\varphi_X \colon \mathsf{Dgm}(X) \to \mathsf{Dgm}(Z)$ and $\varphi_Y \colon  \mathsf{Dgm}(Y) \to \mathsf{Dgm}(Z)$, respectively.
  Define
  \[
    Z' \coloneq \bigoplus_{\substack{(a, b) \in \Dgm(X) \\ (c, d) \in \Dgm(Y) \\ \varphi_X(a, b) = \varphi_Y(c, d)}} K(\max(a, c), \max(b, d)).
  \]
  Observe that $f_b$ and $g_b$ factor through $Z'$. This defines a new span (of strictly bar-to-bar/algebraic matching morphisms) $X \leftarrow Z' \rightarrow Y$ with a smaller cost than the previous span: essentially, for each bar $\langle e, f\rangle$ in $Z$, we consider its matched bars $\langle a, b \rangle$ and $\langle c, d\rangle$ in $X$ and $Y$, respectively, and we move the end-points of $\langle e, f\rangle$ as far left as possible while keeping $f_b$ and $g_b$ well-defined. This makes the cokernels and kernels of these two morphisms smaller.

  By the discussion for the \fbox{$\leq$} case, the cost of this span is equal to the cost of the matching $\psi = \varphi_X \circ \varphi_Y^{-1}$. Thus $\psi \colon \mathsf{Dgm}(X) \to \mathsf{Dgm}(Y)$ is a matching whose cost is smaller than the original span. This concludes the proof.
\end{proof}

\printbibliography